\documentclass[output=paper,colorlinks,citecolor=brown
% ,hidelinks
% showindex
]{langscibook}
\ChapterDOI{10.5281/zenodo.20285196}
%This is where you put the authors and their affiliations
\author{Margaret Matte and Nhu-Anh Nguyen and Emmanuel Bawa and Josh Amaris and Zhilang Liu and Aidan Malanoski and Olivia Mignone and Shane Quinn and Alaa Sharif and Jason Kandybowicz
\affiliation{The Graduate Center, City University of New York}}

%Insert your title here
\title{Polar and \textit{wh-} question formation in Ekhwa Adara}  
\abstract{
This paper presents the first description of polar and \textit{wh-} interrogatives in Ekhwa Adara. In Ekhwa Adara, polar questions are optionally marked by sentence-initial Q particles and involve final-vowel lengthening and L\% boundary tones, while \textit{wh-} question formation involves optional \textit{wh-} movement where the \textit{wh-} item obligatorily precedes a predetermined focus particle. Long-distance \textit{wh-} movement is possible and requires pronominal resumption. Partial \textit{wh-} movement is also possible but seemingly does not require pronominal resumption. Indirect questions are formed via relativization, except when embedded under `ask'. Multiple \textit{wh-} questions exhibit an absence of superiority effects and multiple \textit{wh-} fronting might be available, albeit possibly constrained by an anti-superiority condition.
}

\IfFileExists{../localcommands.tex}{
   \addbibresource{../localbibliography.bib}
   %%%%%%%%%%%%%%%%%%%%%%%%%%%%%%%%%%%%%%%%%%%%%
%%%%%%%%%% From LSP %%%%%%%%%%%%%%%%%%%%%%%%%
%%%%%%%%%%%%%%%%%%%%%%%%%%%%%%%%%%%%%%%%%%%%%

\usepackage{tabularx,multicol}
\usepackage{url}
\urlstyle{same}

\usepackage{langsci-optional}
\usepackage{langsci-lgr}
\usepackage{langsci-gb4e}

% NOTE THAT THE FOLLOWING PACKAGES ARE BANNED: 
% - pstricks
% - tipa
%%%%%%%%%%%%%%%%%%%%%%%%%%%%%%%%%%%%%%%%%%%%%
%%%%%%%%%% From ACAL LaTeX template %%%%%%%%%
%%%%%%%%%%%%%%%%%%%%%%%%%%%%%%%%%%%%%%%%%%%%%

%For stacking text, used here in autosegmental diagrams
\usepackage{stackengine}

%To combine rows in tables
\usepackage{multirow}

%Gives some extra formatting options, e.g. underlining/strikeout
\usepackage[normalem]{ulem}

%Use package/command below to create a double-spaced document, if you want one. Uncomment BOTH the package and the command (\doublespacing) to create a doublespaced document, or leave them as is to have a single-spaced document.
%\usepackage{setspace}
%\doublespacing 

 

%use for special OT tableaux symbols like bomb and sad face. must be loaded early on because it doesn't play well with some other packages
\usepackage{fourier-orns}

%Basic math symbols 
\usepackage{pifont}
\usepackage{amssymb}

%%%Gives shortcuts for glossing. The use of this package is NOT explained in the Quick Reference Guide, but the documentation is on CTAN for those that are interested. MJKD finds it handy for glossing. (https://ctan.org/pkg/leipzig?lang=en)
\usepackage{leipzig}

%Tables
% \usepackage{caption} %For table captions 

%For OT-style tableaux
\usepackage[notipa]{ot-tableau}

%highlights text with \hl{text}
\usepackage{color, soul} %note that soul sometimes eats Unicode text witouth telling you. 

%Drawing Syntax Trees
\usepackage[linguistics]{forest}

%This specifies some formatting for the forest trees to make them nicer to look at. Unfortunately this was borrowed by Michael Diercks from someone a while ago and he can't remember who. Apologies and if someone lets him know he will update this.
\forestset{
  nice nodes/.style={
    for tree={
      inner sep=0pt,
      fit=band,
    },
  },
  default preamble=nice nodes,
}

%A couple of packages that seemed to prefer being called toward the end of the preamble

%This package provides macros for typesetting SPE-style phonological rules. I've redefined the \oneof command here, so that there the rule includes a right brace. 
\usepackage{phonrule}
\renewcommand{\oneof}[2][c]{\ensuremath{
    ~\left\{
    \begin{tabular}{#1#1}#2\end{tabular}
    \right\}
    }}

%For using Greek letters outside of math mode.
% \usepackage{textgreek}


%Random, lets us use the XeLaTeX logo. Not important to the template at all.
% \usepackage{metalogo}



%%%%%%%%%%%%%%%%%%%%%%%%%%%%%%%%%%%%%%%%%%%%%
%%%%%%%%%% From Smith paper %%%%%%%%%%%%%%%%%
%%%%%%%%%%%%%%%%%%%%%%%%%%%%%%%%%%%%%%%%%%%%%

\usetikzlibrary{positioning}
\usetikzlibrary{trees}
% \usepackage{pstricks} %pstricks is banned. Use tikz instead. 
% \usepackage{pst-node}
%\usepackage{tikz}

%%%%%%%%%%%%%%%%%%%%%%%%%%%%%%%%%%%%%%%%%%%%%
%%%%%%%%%% From Gluckman paper %%%%%%%%%%%%%%
%%%%%%%%%%%%%%%%%%%%%%%%%%%%%%%%%%%%%%%%%%%%%

\usepackage{amsmath}



%%%%%%%%%%%%%%%%%%%%%%%%%%%%%%%%%%%%%%%%%%%%%
%%%%%%%%%% From Kandybowicz paper %%%%%%%%%%%
%%%%%%%%%%%%%%%%%%%%%%%%%%%%%%%%%%%%%%%%%%%%%

%These are in the .tex, but there's nothing in the "Figures" folder so I'm not sure that it's actually doing anything. 
 
% \graphicspath{ {./Figures/} }

%%%%%%%%%%%%%%%%%%%%%%%%%%%%%%%%%%%%%%%%%%%%%
%%%%%%%%%% From Matte paper %%%%%%%%%%%%%%%%%
%%%%%%%%%%%%%%%%%%%%%%%%%%%%%%%%%%%%%%%%%%%%%

%%%%%%%%%%%%%%%%%%%%%%%%%%%%%%%%%%%%%%%%%%%%%
%%%%%%%%%% From Grollemund paper %%%%%%%%%%%%%%
%%%%%%%%%%%%%%%%%%%%%%%%%%%%%%%%%%%%%%%%%%%%%

% \usepackage{lscape}

%%%%%%%%%%%%%%%%%%%%%%%%%%%%%%%%%%%%%%%%%%%%%
%%%%%%%%%% From Burns paper %%%%%%%%%%%%%%
%%%%%%%%%%%%%%%%%%%%%%%%%%%%%%%%%%%%%%%%%%%%%


% \usepackage{url}
% \urlstyle{same}

\usepackage{listings}
\lstset{basicstyle=\ttfamily,tabsize=2,breaklines=true}

%MJKD commented out - these are standard for chapters in edited volumes, I don't think we need them here in the preamble. 

% \usepackage{tabularx,multicol}
% \usepackage{langsci-basic}
% \usepackage{langsci-gb4e}
% \bibliography{localbibliography}
% \newcommand{\orcid}[1]{}
% \pagenumbering{roman}

% \IfFileExists{../localcommands.tex}{%hack to check whether this is being compiled as part of a collection or standalone
%    %%%%%%%%%%%%%%%%%%%%%%%%%%%%%%%%%%%%%%%%%%%%%
%%%%%%%%%% From LSP %%%%%%%%%%%%%%%%%%%%%%%%%
%%%%%%%%%%%%%%%%%%%%%%%%%%%%%%%%%%%%%%%%%%%%%

\usepackage{tabularx,multicol}
\usepackage{url}
\urlstyle{same}

\usepackage{langsci-optional}
\usepackage{langsci-lgr}
\usepackage{langsci-gb4e}

% NOTE THAT THE FOLLOWING PACKAGES ARE BANNED: 
% - pstricks
% - tipa
%%%%%%%%%%%%%%%%%%%%%%%%%%%%%%%%%%%%%%%%%%%%%
%%%%%%%%%% From ACAL LaTeX template %%%%%%%%%
%%%%%%%%%%%%%%%%%%%%%%%%%%%%%%%%%%%%%%%%%%%%%

%For stacking text, used here in autosegmental diagrams
\usepackage{stackengine}

%To combine rows in tables
\usepackage{multirow}

%Gives some extra formatting options, e.g. underlining/strikeout
\usepackage[normalem]{ulem}

%Use package/command below to create a double-spaced document, if you want one. Uncomment BOTH the package and the command (\doublespacing) to create a doublespaced document, or leave them as is to have a single-spaced document.
%\usepackage{setspace}
%\doublespacing 

 

%use for special OT tableaux symbols like bomb and sad face. must be loaded early on because it doesn't play well with some other packages
\usepackage{fourier-orns}

%Basic math symbols 
\usepackage{pifont}
\usepackage{amssymb}

%%%Gives shortcuts for glossing. The use of this package is NOT explained in the Quick Reference Guide, but the documentation is on CTAN for those that are interested. MJKD finds it handy for glossing. (https://ctan.org/pkg/leipzig?lang=en)
\usepackage{leipzig}

%Tables
% \usepackage{caption} %For table captions 

%For OT-style tableaux
\usepackage[notipa]{ot-tableau}

%highlights text with \hl{text}
\usepackage{color, soul} %note that soul sometimes eats Unicode text witouth telling you. 

%Drawing Syntax Trees
\usepackage[linguistics]{forest}

%This specifies some formatting for the forest trees to make them nicer to look at. Unfortunately this was borrowed by Michael Diercks from someone a while ago and he can't remember who. Apologies and if someone lets him know he will update this.
\forestset{
  nice nodes/.style={
    for tree={
      inner sep=0pt,
      fit=band,
    },
  },
  default preamble=nice nodes,
}

%A couple of packages that seemed to prefer being called toward the end of the preamble

%This package provides macros for typesetting SPE-style phonological rules. I've redefined the \oneof command here, so that there the rule includes a right brace. 
\usepackage{phonrule}
\renewcommand{\oneof}[2][c]{\ensuremath{
    ~\left\{
    \begin{tabular}{#1#1}#2\end{tabular}
    \right\}
    }}

%For using Greek letters outside of math mode.
% \usepackage{textgreek}


%Random, lets us use the XeLaTeX logo. Not important to the template at all.
% \usepackage{metalogo}



%%%%%%%%%%%%%%%%%%%%%%%%%%%%%%%%%%%%%%%%%%%%%
%%%%%%%%%% From Smith paper %%%%%%%%%%%%%%%%%
%%%%%%%%%%%%%%%%%%%%%%%%%%%%%%%%%%%%%%%%%%%%%

\usetikzlibrary{positioning}
\usetikzlibrary{trees}
% \usepackage{pstricks} %pstricks is banned. Use tikz instead. 
% \usepackage{pst-node}
%\usepackage{tikz}

%%%%%%%%%%%%%%%%%%%%%%%%%%%%%%%%%%%%%%%%%%%%%
%%%%%%%%%% From Gluckman paper %%%%%%%%%%%%%%
%%%%%%%%%%%%%%%%%%%%%%%%%%%%%%%%%%%%%%%%%%%%%

\usepackage{amsmath}



%%%%%%%%%%%%%%%%%%%%%%%%%%%%%%%%%%%%%%%%%%%%%
%%%%%%%%%% From Kandybowicz paper %%%%%%%%%%%
%%%%%%%%%%%%%%%%%%%%%%%%%%%%%%%%%%%%%%%%%%%%%

%These are in the .tex, but there's nothing in the "Figures" folder so I'm not sure that it's actually doing anything. 
 
% \graphicspath{ {./Figures/} }

%%%%%%%%%%%%%%%%%%%%%%%%%%%%%%%%%%%%%%%%%%%%%
%%%%%%%%%% From Matte paper %%%%%%%%%%%%%%%%%
%%%%%%%%%%%%%%%%%%%%%%%%%%%%%%%%%%%%%%%%%%%%%

%%%%%%%%%%%%%%%%%%%%%%%%%%%%%%%%%%%%%%%%%%%%%
%%%%%%%%%% From Grollemund paper %%%%%%%%%%%%%%
%%%%%%%%%%%%%%%%%%%%%%%%%%%%%%%%%%%%%%%%%%%%%

% \usepackage{lscape}

%%%%%%%%%%%%%%%%%%%%%%%%%%%%%%%%%%%%%%%%%%%%%
%%%%%%%%%% From Burns paper %%%%%%%%%%%%%%
%%%%%%%%%%%%%%%%%%%%%%%%%%%%%%%%%%%%%%%%%%%%%


% \usepackage{url}
% \urlstyle{same}

\usepackage{listings}
\lstset{basicstyle=\ttfamily,tabsize=2,breaklines=true}

%MJKD commented out - these are standard for chapters in edited volumes, I don't think we need them here in the preamble. 

% \usepackage{tabularx,multicol}
% \usepackage{langsci-basic}
% \usepackage{langsci-gb4e}
% \bibliography{localbibliography}
% \newcommand{\orcid}[1]{}
% \pagenumbering{roman}

% \IfFileExists{../localcommands.tex}{%hack to check whether this is being compiled as part of a collection or standalone
%    %%%%%%%%%%%%%%%%%%%%%%%%%%%%%%%%%%%%%%%%%%%%%
%%%%%%%%%% From LSP %%%%%%%%%%%%%%%%%%%%%%%%%
%%%%%%%%%%%%%%%%%%%%%%%%%%%%%%%%%%%%%%%%%%%%%

\usepackage{tabularx,multicol}
\usepackage{url}
\urlstyle{same}

\usepackage{langsci-optional}
\usepackage{langsci-lgr}
\usepackage{langsci-gb4e}

% NOTE THAT THE FOLLOWING PACKAGES ARE BANNED: 
% - pstricks
% - tipa
%%%%%%%%%%%%%%%%%%%%%%%%%%%%%%%%%%%%%%%%%%%%%
%%%%%%%%%% From ACAL LaTeX template %%%%%%%%%
%%%%%%%%%%%%%%%%%%%%%%%%%%%%%%%%%%%%%%%%%%%%%

%For stacking text, used here in autosegmental diagrams
\usepackage{stackengine}

%To combine rows in tables
\usepackage{multirow}

%Gives some extra formatting options, e.g. underlining/strikeout
\usepackage[normalem]{ulem}

%Use package/command below to create a double-spaced document, if you want one. Uncomment BOTH the package and the command (\doublespacing) to create a doublespaced document, or leave them as is to have a single-spaced document.
%\usepackage{setspace}
%\doublespacing 

 

%use for special OT tableaux symbols like bomb and sad face. must be loaded early on because it doesn't play well with some other packages
\usepackage{fourier-orns}

%Basic math symbols 
\usepackage{pifont}
\usepackage{amssymb}

%%%Gives shortcuts for glossing. The use of this package is NOT explained in the Quick Reference Guide, but the documentation is on CTAN for those that are interested. MJKD finds it handy for glossing. (https://ctan.org/pkg/leipzig?lang=en)
\usepackage{leipzig}

%Tables
% \usepackage{caption} %For table captions 

%For OT-style tableaux
\usepackage[notipa]{ot-tableau}

%highlights text with \hl{text}
\usepackage{color, soul} %note that soul sometimes eats Unicode text witouth telling you. 

%Drawing Syntax Trees
\usepackage[linguistics]{forest}

%This specifies some formatting for the forest trees to make them nicer to look at. Unfortunately this was borrowed by Michael Diercks from someone a while ago and he can't remember who. Apologies and if someone lets him know he will update this.
\forestset{
  nice nodes/.style={
    for tree={
      inner sep=0pt,
      fit=band,
    },
  },
  default preamble=nice nodes,
}

%A couple of packages that seemed to prefer being called toward the end of the preamble

%This package provides macros for typesetting SPE-style phonological rules. I've redefined the \oneof command here, so that there the rule includes a right brace. 
\usepackage{phonrule}
\renewcommand{\oneof}[2][c]{\ensuremath{
    ~\left\{
    \begin{tabular}{#1#1}#2\end{tabular}
    \right\}
    }}

%For using Greek letters outside of math mode.
% \usepackage{textgreek}


%Random, lets us use the XeLaTeX logo. Not important to the template at all.
% \usepackage{metalogo}



%%%%%%%%%%%%%%%%%%%%%%%%%%%%%%%%%%%%%%%%%%%%%
%%%%%%%%%% From Smith paper %%%%%%%%%%%%%%%%%
%%%%%%%%%%%%%%%%%%%%%%%%%%%%%%%%%%%%%%%%%%%%%

\usetikzlibrary{positioning}
\usetikzlibrary{trees}
% \usepackage{pstricks} %pstricks is banned. Use tikz instead. 
% \usepackage{pst-node}
%\usepackage{tikz}

%%%%%%%%%%%%%%%%%%%%%%%%%%%%%%%%%%%%%%%%%%%%%
%%%%%%%%%% From Gluckman paper %%%%%%%%%%%%%%
%%%%%%%%%%%%%%%%%%%%%%%%%%%%%%%%%%%%%%%%%%%%%

\usepackage{amsmath}



%%%%%%%%%%%%%%%%%%%%%%%%%%%%%%%%%%%%%%%%%%%%%
%%%%%%%%%% From Kandybowicz paper %%%%%%%%%%%
%%%%%%%%%%%%%%%%%%%%%%%%%%%%%%%%%%%%%%%%%%%%%

%These are in the .tex, but there's nothing in the "Figures" folder so I'm not sure that it's actually doing anything. 
 
% \graphicspath{ {./Figures/} }

%%%%%%%%%%%%%%%%%%%%%%%%%%%%%%%%%%%%%%%%%%%%%
%%%%%%%%%% From Matte paper %%%%%%%%%%%%%%%%%
%%%%%%%%%%%%%%%%%%%%%%%%%%%%%%%%%%%%%%%%%%%%%

%%%%%%%%%%%%%%%%%%%%%%%%%%%%%%%%%%%%%%%%%%%%%
%%%%%%%%%% From Grollemund paper %%%%%%%%%%%%%%
%%%%%%%%%%%%%%%%%%%%%%%%%%%%%%%%%%%%%%%%%%%%%

% \usepackage{lscape}

%%%%%%%%%%%%%%%%%%%%%%%%%%%%%%%%%%%%%%%%%%%%%
%%%%%%%%%% From Burns paper %%%%%%%%%%%%%%
%%%%%%%%%%%%%%%%%%%%%%%%%%%%%%%%%%%%%%%%%%%%%


% \usepackage{url}
% \urlstyle{same}

\usepackage{listings}
\lstset{basicstyle=\ttfamily,tabsize=2,breaklines=true}

%MJKD commented out - these are standard for chapters in edited volumes, I don't think we need them here in the preamble. 

% \usepackage{tabularx,multicol}
% \usepackage{langsci-basic}
% \usepackage{langsci-gb4e}
% \bibliography{localbibliography}
% \newcommand{\orcid}[1]{}
% \pagenumbering{roman}

% \IfFileExists{../localcommands.tex}{%hack to check whether this is being compiled as part of a collection or standalone
%    \input{../localpackages}
%    \input{../localcommands}
%    \input{../localhyphenation}
%     \bibliography{localbibliography}
%     \togglepaper[23]
% }{}

% %custom footer for preprints
% \papernote{\scriptsize\normalfont
%     Change author.
%     Change title.
%     To appear in:
%     Change Volume Editor.  
%     Change volume title.  
%     Berlin: Language Science Press. [preliminary page numbering]
% }



%%%%%%%%%%%%%%%%%%%%%%%%%%%%%%%%%%%%%%%%%%%%%
%%%%%%%%%% From Feibig paper %%%%%%%%%%%%%%
%%%%%%%%%%%%%%%%%%%%%%%%%%%%%%%%%%%%%%%%%%%%%

\usepackage{tikz-qtree}
\usepackage{tikz-qtree-compat}

%%%%%%%%%%%%%%%%%%%%%%%%%%%%%%%%%%%%%%%%%%%%%
%%%%%%%%%% From Olson paper %%%%%%%%%%%%%%
%%%%%%%%%%%%%%%%%%%%%%%%%%%%%%%%%%%%%%%%%%%%%

%MJKD - TIPA is forbidden lol. Commenting out for now but we'll have to address this in the Olson paper

%\usepackage[tone]{tipa}

%%%%%%%%%%%%%%%%%%%%%%%%%%%%%%%%%%%%%%%%%%%%%
%%%%%%%%%% From Caragine paper %%%%%%%%%%%%%%
%%%%%%%%%%%%%%%%%%%%%%%%%%%%%%%%%%%%%%%%%%%%%

\usepackage{placeins}
% \usepackage{graphicx}
% \usepackage{float} %Overleaf suggested this as a solution


%%%%%%%%%%%%%%%%%%%%%%%%%%%%%%%%%%%%%%%%%%%%%
%%%%%%%%%% From Kieffer paper %%%%%%%%%%%%%%
%%%%%%%%%%%%%%%%%%%%%%%%%%%%%%%%%%%%%%%%%%%%%

\usepackage{array}

%%%%%%%%%%%%%%%%%%%%%%%%%%%%%%%%%%%%%%%%%%%%%
%%%%%%%%%% From Payne paper %%%%%%%%%%%%%%
%%%%%%%%%%%%%%%%%%%%%%%%%%%%%%%%%%%%%%%%%%%%%

\usepackage{enumerate}
\usepackage{array,ragged2e}
\usepackage{pgfplots}


%%%%%%%%%%%%%%%%%%%%%%%%%%%%%%%%%%%%%%%%%%%%%
%%%%%%%%%% From Diercks paper %%%%%%%%%%%%%%
%%%%%%%%%%%%%%%%%%%%%%%%%%%%%%%%%%%%%%%%%%%%%

% \usepackage{cgloss}



%%%%%%%%%%%%%%%%%%%%%%%%%%%%%%%%%%%%%%%%%%%%%
%%%%%%%%%% From Li paper %%%%%%%%%%%%%%
%%%%%%%%%%%%%%%%%%%%%%%%%%%%%%%%%%%%%%%%%%%%%



%%%%%%%%%%%%%%%%%%%%%%%%%%%%%%%%%%%%%%%%%%%%%
%%%%%%%%%% From Myers paper %%%%%%%%%%%%%%
%%%%%%%%%%%%%%%%%%%%%%%%%%%%%%%%%%%%%%%%%%%%%



\usepackage{tikz-qtree}
\usepackage{tikz-qtree-compat}


%Package that makes glossing slightly easier.
\RequirePackage{leipzig}
%\RequirePackage{mathtools} % for \mathrlap

% % % Shortcuts (various borrowed from Jason Zentz)
\RequirePackage{xspace}
\xspaceaddexceptions{]\}}
\usepackage{langsci-textipa}
\usepackage{langsci-branding}
\usepackage{tabto}

%    %%%%%%%%%%%%%%%%%%%%%%%%%%%%%%%%%%%%%%%%%%%%%
%%%%%%%%%% From LSP %%%%%%%%%%%%%%%%%%%%%%%%%
%%%%%%%%%%%%%%%%%%%%%%%%%%%%%%%%%%%%%%%%%%%%%


% \newcommand*{\orcid}{}


\makeatletter
\let\thetitle\@title
\let\theauthor\@author
\makeatother

\newcommand{\togglepaper}[1][0]{
   \bibliography{../localbibliography}
   \papernote{\scriptsize\normalfont
     \theauthor.
     \titleTemp.
     To appear in:
     E. Di Tor \& Herr Rausgeberin (ed.).
     Booktitle in localcommands.tex.
     Berlin: Language Science Press. [preliminary page numbering]
   }
   \pagenumbering{roman}
   \setcounter{chapter}{#1}
   \addtocounter{chapter}{-1}
}

\newbool{bookcompile}
\booltrue{bookcompile}
\newcommand{\bookorchapter}[2]{\ifbool{bookcompile}{#1}{#2}}


%%%%%%%%%%%%%%%%%%%%%%%%%%%%%%%%%%%%%%%%%%%%%
%%%%%%%%%% From ACAL LaTeX template %%%%%%%%%
%%%%%%%%%%%%%%%%%%%%%%%%%%%%%%%%%%%%%%%%%%%%%


\usepackage{tikz}

\newcommand{\circled}[1]{\begin{tikzpicture}[baseline=(word.base)]
\node[draw, rounded corners, text height=8pt, text depth=2pt, inner sep=2pt, outer sep=0pt, use as bounding box] (word) {#1};
\end{tikzpicture}
}


%%%%%%%%%%%%%%%%%%%%%%%%%%%%%%%%%%%%%%%%%%%%%%
%%%%%%%%%%%%%%%%%%%%%%%%%%%%%%%%%%%%%%%%%%%%%%

% Useful Ling Shortcuts


%This makes the \emptyset command be a nicer one
\let\oldemptyset\emptyset
\let\emptyset\varnothing
\newcommand{\nothing}{$\emptyset$}

%Not all of these are explained in the Quick Reference Guide, but they are here bc they are relevant to some of our students. 
\newcommand{\xbar}{\ensuremath{'}\xspace}
\newcommand{\xhead}{\ensuremath{^\circ}\xspace}
\newcommand{\Lb}[1]{$\text{[}_{\text{#1}}$ } %A more convenient left bracket
\newcommand{\Rb}[1]{$\text{]}_{\text{#1}}$ } %A more convenient left bracket
\newcommand{\gap}{\underline{\hspace{1.2em}}}
\newcommand{\vP}{\emph{v}P}
\newcommand{\lilv}{\emph{v}}
\newcommand{\Abar}{A$'$-} %A more convenient A-bar notation
\newcommand{\ph}{$\varphi$\xspace} %A more convenient phi
\newcommand{\pro}{\emph{pro}\xspace}
\newcommand{\subs}[1]{\textsubscript{#1}} %A more convenient subscript
\newcommand{\spells}{$\Longleftrightarrow$} %spellout arrow for morph spellout rules
\newcommand{\tr}[1]{\textit{t}\textsubscript{\textit{#1}}} %easy traces with subscript
\newcommand{\supers}[1]{\textsuperscript{#1}}

%These are borrowed/modified from Andrew Mackenzie's ling-macros package. We have copied them in here (instead of just using the package itself) to avoid a minor clash that was generating an error.


% Abbreviations for glossing, based on Leipzig
\newleipzig{hab}{hab}{habitual}
\newleipzig{rem}{rem}{remote}
\newleipzig{sm}{sm}{subject marker}
\newleipzig{t}{t}{tense}
\newleipzig{aa}{aa}{anti-agreement}
\newleipzig{pron}{pron}{pronoun}
\newleipzig{rec}{rec}{recent}
\newleipzig{om}{om}{object marker}
%\newleipzig{ipfv}{ipfv}{imperfective}
\newleipzig{asp}{asp}{aspect}
\newleipzig{lk}{lk}{linker}
\newleipzig{pcl}{pcl}{particle}
\newleipzig{stat}{stat}{stative}
\newleipzig{ints}{ints}{intensive}
\newleipzig{ascl}{ascl}{assertive subject clitic}
\newleipzig{nascl}{nascl}{non-assertive subject clitic}
\newleipzig{ta}{ta}{tense and/or aspect}
\newleipzig{assoc}{assoc}{associative marker}
\newleipzig{hon}{hon}{honorific}
%\newleipzig{whprt}{wh}{\wh particle}
\newleipzig{sa}{sa}{subject agreement}
\newleipzig{conj}{conj}{conjunction}
%\newleipzig{loc}{loc}{locative}
\newleipzig{expl}{expl}{expletive}
\newleipzig{rcm}{rcm}{reciprocal marker}
\newleipzig{pers}{pers}{persistive}
\newleipzig{fv}{fv}{final vowel}
%\newleipzig{}{}{} %this is just to copy for when I want to add more


%%%%%%%%%%%%%%%%%%%%%%%%%%%%%%%%%%%%%%%%%%%%%%%%%%%%%
%%%%%%%%%%%%%%%%%%%%%%%%%%%%%%%%%%%%%%%%%%%%%%%%%%%%%


%%%%%%%%%%%%%%%%%%%%%%%%%%%%%%%%%%%%%%%%%%%%%
%%%%%%%%%% From Gluckman paper %%%%%%%%%%%%%%
%%%%%%%%%%%%%%%%%%%%%%%%%%%%%%%%%%%%%%%%%%%%%

\newcommand{\pcite}[1]{\citeauthor{#1}'s (\citeyear{#1})}


%%%%%%%%%%%%%%%%%%%%%%%%%%%%%%%%%%%%%%%%%%%%%
%%%%%%%%%% From Myers paper %%%%%%%%%%%%%%
%%%%%%%%%%%%%%%%%%%%%%%%%%%%%%%%%%%%%%%%%%%%%

%% hspace over width of something without showing it
\newcommand*{\hspaceThis}[1]{\settowidth{\LSPTmp}{#1}\hspace*{\LSPTmp}}
% use \hphantom{} for this


%%%%%%%%%%%%%%%%%%%%%%%%%%%%%%%%%%%%%%%%%%%%%
%%%%%%%%%% From Matte paper %%%%%%%%%%%%%%%%%
%%%%%%%%%%%%%%%%%%%%%%%%%%%%%%%%%%%%%%%%%%%%%

%mjkd commented this out in case it's breaking something right now. 
% \newcounter{xlistex}
% \newcommand{\footnoteex}{\stepcounter{xlistex}(\roman{xlistex})}

\newleipzig{sp}{sp}{speaker} %a new leipzig glossing command


%%%%%%%%%%%%%%%%%%%%%%%%%%%%%%%%%%%%%%%%%%%%%
%%%%%%%%%% From GamFranich paper %%%%%%%%%%%%%%
%%%%%%%%%%%%%%%%%%%%%%%%%%%%%%%%%%%%%%%%%%%%%

%MJKD commented these out - both will be provided by the overall LSP book template 

% \bibliography{localbibliography}
% \newcommand{\orcid}[1]{}

%%%%%%%%%%%%%%%%%%%%%%%%%%%%%%%%%%%%%%%%%%%%%
%%%%%%%%%% From Diercks paper %%%%%%%%%%%%%%
%%%%%%%%%%%%%%%%%%%%%%%%%%%%%%%%%%%%%%%%%%%%%

\newcommand{\abar}{$\bar{\text{A}}$\xspace} % this one is different to the \Abar command in Mike's shortcuts doc
\newcommand{\topFeat}{[\textsc{top}]\xspace}


%%%%%%%%%%%%%%%%%%%%%%%%%%%%%%%%%%%%%%%%%%%%%
%%%%%%%%%% From Kinchsular paper %%%%%%%%%%%%%%
%%%%%%%%%%%%%%%%%%%%%%%%%%%%%%%%%%%%%%%%%%%%%

%%%%%%%%%%%%%%%%%%%%%%%%%%%%%%%%%%%%%%%%%%%%%
%%%%%%%%%% From Chen paper %%%%%%%%%%%%%%
%%%%%%%%%%%%%%%%%%%%%%%%%%%%%%%%%%%%%%%%%%%%%

\newcounter{savedex}
\makeatletter
\newcommand*{\saveex}{\setcounter{savedex}{\value{\@xnumctr}}}
\newcommand*{\resumeex}{\setcounter{\@xnumctr}{\value{savedex}}}
\makeatother


%    \input{../localhyphenation}
%     \bibliography{localbibliography}
%     \togglepaper[23]
% }{}

% %custom footer for preprints
% \papernote{\scriptsize\normalfont
%     Change author.
%     Change title.
%     To appear in:
%     Change Volume Editor.  
%     Change volume title.  
%     Berlin: Language Science Press. [preliminary page numbering]
% }



%%%%%%%%%%%%%%%%%%%%%%%%%%%%%%%%%%%%%%%%%%%%%
%%%%%%%%%% From Feibig paper %%%%%%%%%%%%%%
%%%%%%%%%%%%%%%%%%%%%%%%%%%%%%%%%%%%%%%%%%%%%

\usepackage{tikz-qtree}
\usepackage{tikz-qtree-compat}

%%%%%%%%%%%%%%%%%%%%%%%%%%%%%%%%%%%%%%%%%%%%%
%%%%%%%%%% From Olson paper %%%%%%%%%%%%%%
%%%%%%%%%%%%%%%%%%%%%%%%%%%%%%%%%%%%%%%%%%%%%

%MJKD - TIPA is forbidden lol. Commenting out for now but we'll have to address this in the Olson paper

%\usepackage[tone]{tipa}

%%%%%%%%%%%%%%%%%%%%%%%%%%%%%%%%%%%%%%%%%%%%%
%%%%%%%%%% From Caragine paper %%%%%%%%%%%%%%
%%%%%%%%%%%%%%%%%%%%%%%%%%%%%%%%%%%%%%%%%%%%%

\usepackage{placeins}
% \usepackage{graphicx}
% \usepackage{float} %Overleaf suggested this as a solution


%%%%%%%%%%%%%%%%%%%%%%%%%%%%%%%%%%%%%%%%%%%%%
%%%%%%%%%% From Kieffer paper %%%%%%%%%%%%%%
%%%%%%%%%%%%%%%%%%%%%%%%%%%%%%%%%%%%%%%%%%%%%

\usepackage{array}

%%%%%%%%%%%%%%%%%%%%%%%%%%%%%%%%%%%%%%%%%%%%%
%%%%%%%%%% From Payne paper %%%%%%%%%%%%%%
%%%%%%%%%%%%%%%%%%%%%%%%%%%%%%%%%%%%%%%%%%%%%

\usepackage{enumerate}
\usepackage{array,ragged2e}
\usepackage{pgfplots}


%%%%%%%%%%%%%%%%%%%%%%%%%%%%%%%%%%%%%%%%%%%%%
%%%%%%%%%% From Diercks paper %%%%%%%%%%%%%%
%%%%%%%%%%%%%%%%%%%%%%%%%%%%%%%%%%%%%%%%%%%%%

% \usepackage{cgloss}



%%%%%%%%%%%%%%%%%%%%%%%%%%%%%%%%%%%%%%%%%%%%%
%%%%%%%%%% From Li paper %%%%%%%%%%%%%%
%%%%%%%%%%%%%%%%%%%%%%%%%%%%%%%%%%%%%%%%%%%%%



%%%%%%%%%%%%%%%%%%%%%%%%%%%%%%%%%%%%%%%%%%%%%
%%%%%%%%%% From Myers paper %%%%%%%%%%%%%%
%%%%%%%%%%%%%%%%%%%%%%%%%%%%%%%%%%%%%%%%%%%%%



\usepackage{tikz-qtree}
\usepackage{tikz-qtree-compat}


%Package that makes glossing slightly easier.
\RequirePackage{leipzig}
%\RequirePackage{mathtools} % for \mathrlap

% % % Shortcuts (various borrowed from Jason Zentz)
\RequirePackage{xspace}
\xspaceaddexceptions{]\}}
\usepackage{langsci-textipa}
\usepackage{langsci-branding}
\usepackage{tabto}

%    %%%%%%%%%%%%%%%%%%%%%%%%%%%%%%%%%%%%%%%%%%%%%
%%%%%%%%%% From LSP %%%%%%%%%%%%%%%%%%%%%%%%%
%%%%%%%%%%%%%%%%%%%%%%%%%%%%%%%%%%%%%%%%%%%%%


% \newcommand*{\orcid}{}


\makeatletter
\let\thetitle\@title
\let\theauthor\@author
\makeatother

\newcommand{\togglepaper}[1][0]{
   \bibliography{../localbibliography}
   \papernote{\scriptsize\normalfont
     \theauthor.
     \titleTemp.
     To appear in:
     E. Di Tor \& Herr Rausgeberin (ed.).
     Booktitle in localcommands.tex.
     Berlin: Language Science Press. [preliminary page numbering]
   }
   \pagenumbering{roman}
   \setcounter{chapter}{#1}
   \addtocounter{chapter}{-1}
}

\newbool{bookcompile}
\booltrue{bookcompile}
\newcommand{\bookorchapter}[2]{\ifbool{bookcompile}{#1}{#2}}


%%%%%%%%%%%%%%%%%%%%%%%%%%%%%%%%%%%%%%%%%%%%%
%%%%%%%%%% From ACAL LaTeX template %%%%%%%%%
%%%%%%%%%%%%%%%%%%%%%%%%%%%%%%%%%%%%%%%%%%%%%


\usepackage{tikz}

\newcommand{\circled}[1]{\begin{tikzpicture}[baseline=(word.base)]
\node[draw, rounded corners, text height=8pt, text depth=2pt, inner sep=2pt, outer sep=0pt, use as bounding box] (word) {#1};
\end{tikzpicture}
}


%%%%%%%%%%%%%%%%%%%%%%%%%%%%%%%%%%%%%%%%%%%%%%
%%%%%%%%%%%%%%%%%%%%%%%%%%%%%%%%%%%%%%%%%%%%%%

% Useful Ling Shortcuts


%This makes the \emptyset command be a nicer one
\let\oldemptyset\emptyset
\let\emptyset\varnothing
\newcommand{\nothing}{$\emptyset$}

%Not all of these are explained in the Quick Reference Guide, but they are here bc they are relevant to some of our students. 
\newcommand{\xbar}{\ensuremath{'}\xspace}
\newcommand{\xhead}{\ensuremath{^\circ}\xspace}
\newcommand{\Lb}[1]{$\text{[}_{\text{#1}}$ } %A more convenient left bracket
\newcommand{\Rb}[1]{$\text{]}_{\text{#1}}$ } %A more convenient left bracket
\newcommand{\gap}{\underline{\hspace{1.2em}}}
\newcommand{\vP}{\emph{v}P}
\newcommand{\lilv}{\emph{v}}
\newcommand{\Abar}{A$'$-} %A more convenient A-bar notation
\newcommand{\ph}{$\varphi$\xspace} %A more convenient phi
\newcommand{\pro}{\emph{pro}\xspace}
\newcommand{\subs}[1]{\textsubscript{#1}} %A more convenient subscript
\newcommand{\spells}{$\Longleftrightarrow$} %spellout arrow for morph spellout rules
\newcommand{\tr}[1]{\textit{t}\textsubscript{\textit{#1}}} %easy traces with subscript
\newcommand{\supers}[1]{\textsuperscript{#1}}

%These are borrowed/modified from Andrew Mackenzie's ling-macros package. We have copied them in here (instead of just using the package itself) to avoid a minor clash that was generating an error.


% Abbreviations for glossing, based on Leipzig
\newleipzig{hab}{hab}{habitual}
\newleipzig{rem}{rem}{remote}
\newleipzig{sm}{sm}{subject marker}
\newleipzig{t}{t}{tense}
\newleipzig{aa}{aa}{anti-agreement}
\newleipzig{pron}{pron}{pronoun}
\newleipzig{rec}{rec}{recent}
\newleipzig{om}{om}{object marker}
%\newleipzig{ipfv}{ipfv}{imperfective}
\newleipzig{asp}{asp}{aspect}
\newleipzig{lk}{lk}{linker}
\newleipzig{pcl}{pcl}{particle}
\newleipzig{stat}{stat}{stative}
\newleipzig{ints}{ints}{intensive}
\newleipzig{ascl}{ascl}{assertive subject clitic}
\newleipzig{nascl}{nascl}{non-assertive subject clitic}
\newleipzig{ta}{ta}{tense and/or aspect}
\newleipzig{assoc}{assoc}{associative marker}
\newleipzig{hon}{hon}{honorific}
%\newleipzig{whprt}{wh}{\wh particle}
\newleipzig{sa}{sa}{subject agreement}
\newleipzig{conj}{conj}{conjunction}
%\newleipzig{loc}{loc}{locative}
\newleipzig{expl}{expl}{expletive}
\newleipzig{rcm}{rcm}{reciprocal marker}
\newleipzig{pers}{pers}{persistive}
\newleipzig{fv}{fv}{final vowel}
%\newleipzig{}{}{} %this is just to copy for when I want to add more


%%%%%%%%%%%%%%%%%%%%%%%%%%%%%%%%%%%%%%%%%%%%%%%%%%%%%
%%%%%%%%%%%%%%%%%%%%%%%%%%%%%%%%%%%%%%%%%%%%%%%%%%%%%


%%%%%%%%%%%%%%%%%%%%%%%%%%%%%%%%%%%%%%%%%%%%%
%%%%%%%%%% From Gluckman paper %%%%%%%%%%%%%%
%%%%%%%%%%%%%%%%%%%%%%%%%%%%%%%%%%%%%%%%%%%%%

\newcommand{\pcite}[1]{\citeauthor{#1}'s (\citeyear{#1})}


%%%%%%%%%%%%%%%%%%%%%%%%%%%%%%%%%%%%%%%%%%%%%
%%%%%%%%%% From Myers paper %%%%%%%%%%%%%%
%%%%%%%%%%%%%%%%%%%%%%%%%%%%%%%%%%%%%%%%%%%%%

%% hspace over width of something without showing it
\newcommand*{\hspaceThis}[1]{\settowidth{\LSPTmp}{#1}\hspace*{\LSPTmp}}
% use \hphantom{} for this


%%%%%%%%%%%%%%%%%%%%%%%%%%%%%%%%%%%%%%%%%%%%%
%%%%%%%%%% From Matte paper %%%%%%%%%%%%%%%%%
%%%%%%%%%%%%%%%%%%%%%%%%%%%%%%%%%%%%%%%%%%%%%

%mjkd commented this out in case it's breaking something right now. 
% \newcounter{xlistex}
% \newcommand{\footnoteex}{\stepcounter{xlistex}(\roman{xlistex})}

\newleipzig{sp}{sp}{speaker} %a new leipzig glossing command


%%%%%%%%%%%%%%%%%%%%%%%%%%%%%%%%%%%%%%%%%%%%%
%%%%%%%%%% From GamFranich paper %%%%%%%%%%%%%%
%%%%%%%%%%%%%%%%%%%%%%%%%%%%%%%%%%%%%%%%%%%%%

%MJKD commented these out - both will be provided by the overall LSP book template 

% \bibliography{localbibliography}
% \newcommand{\orcid}[1]{}

%%%%%%%%%%%%%%%%%%%%%%%%%%%%%%%%%%%%%%%%%%%%%
%%%%%%%%%% From Diercks paper %%%%%%%%%%%%%%
%%%%%%%%%%%%%%%%%%%%%%%%%%%%%%%%%%%%%%%%%%%%%

\newcommand{\abar}{$\bar{\text{A}}$\xspace} % this one is different to the \Abar command in Mike's shortcuts doc
\newcommand{\topFeat}{[\textsc{top}]\xspace}


%%%%%%%%%%%%%%%%%%%%%%%%%%%%%%%%%%%%%%%%%%%%%
%%%%%%%%%% From Kinchsular paper %%%%%%%%%%%%%%
%%%%%%%%%%%%%%%%%%%%%%%%%%%%%%%%%%%%%%%%%%%%%

%%%%%%%%%%%%%%%%%%%%%%%%%%%%%%%%%%%%%%%%%%%%%
%%%%%%%%%% From Chen paper %%%%%%%%%%%%%%
%%%%%%%%%%%%%%%%%%%%%%%%%%%%%%%%%%%%%%%%%%%%%

\newcounter{savedex}
\makeatletter
\newcommand*{\saveex}{\setcounter{savedex}{\value{\@xnumctr}}}
\newcommand*{\resumeex}{\setcounter{\@xnumctr}{\value{savedex}}}
\makeatother


%    \input{../localhyphenation}
%     \bibliography{localbibliography}
%     \togglepaper[23]
% }{}

% %custom footer for preprints
% \papernote{\scriptsize\normalfont
%     Change author.
%     Change title.
%     To appear in:
%     Change Volume Editor.  
%     Change volume title.  
%     Berlin: Language Science Press. [preliminary page numbering]
% }



%%%%%%%%%%%%%%%%%%%%%%%%%%%%%%%%%%%%%%%%%%%%%
%%%%%%%%%% From Feibig paper %%%%%%%%%%%%%%
%%%%%%%%%%%%%%%%%%%%%%%%%%%%%%%%%%%%%%%%%%%%%

\usepackage{tikz-qtree}
\usepackage{tikz-qtree-compat}

%%%%%%%%%%%%%%%%%%%%%%%%%%%%%%%%%%%%%%%%%%%%%
%%%%%%%%%% From Olson paper %%%%%%%%%%%%%%
%%%%%%%%%%%%%%%%%%%%%%%%%%%%%%%%%%%%%%%%%%%%%

%MJKD - TIPA is forbidden lol. Commenting out for now but we'll have to address this in the Olson paper

%\usepackage[tone]{tipa}

%%%%%%%%%%%%%%%%%%%%%%%%%%%%%%%%%%%%%%%%%%%%%
%%%%%%%%%% From Caragine paper %%%%%%%%%%%%%%
%%%%%%%%%%%%%%%%%%%%%%%%%%%%%%%%%%%%%%%%%%%%%

\usepackage{placeins}
% \usepackage{graphicx}
% \usepackage{float} %Overleaf suggested this as a solution


%%%%%%%%%%%%%%%%%%%%%%%%%%%%%%%%%%%%%%%%%%%%%
%%%%%%%%%% From Kieffer paper %%%%%%%%%%%%%%
%%%%%%%%%%%%%%%%%%%%%%%%%%%%%%%%%%%%%%%%%%%%%

\usepackage{array}

%%%%%%%%%%%%%%%%%%%%%%%%%%%%%%%%%%%%%%%%%%%%%
%%%%%%%%%% From Payne paper %%%%%%%%%%%%%%
%%%%%%%%%%%%%%%%%%%%%%%%%%%%%%%%%%%%%%%%%%%%%

\usepackage{enumerate}
\usepackage{array,ragged2e}
\usepackage{pgfplots}


%%%%%%%%%%%%%%%%%%%%%%%%%%%%%%%%%%%%%%%%%%%%%
%%%%%%%%%% From Diercks paper %%%%%%%%%%%%%%
%%%%%%%%%%%%%%%%%%%%%%%%%%%%%%%%%%%%%%%%%%%%%

% \usepackage{cgloss}



%%%%%%%%%%%%%%%%%%%%%%%%%%%%%%%%%%%%%%%%%%%%%
%%%%%%%%%% From Li paper %%%%%%%%%%%%%%
%%%%%%%%%%%%%%%%%%%%%%%%%%%%%%%%%%%%%%%%%%%%%



%%%%%%%%%%%%%%%%%%%%%%%%%%%%%%%%%%%%%%%%%%%%%
%%%%%%%%%% From Myers paper %%%%%%%%%%%%%%
%%%%%%%%%%%%%%%%%%%%%%%%%%%%%%%%%%%%%%%%%%%%%



\usepackage{tikz-qtree}
\usepackage{tikz-qtree-compat}


%Package that makes glossing slightly easier.
\RequirePackage{leipzig}
%\RequirePackage{mathtools} % for \mathrlap

% % % Shortcuts (various borrowed from Jason Zentz)
\RequirePackage{xspace}
\xspaceaddexceptions{]\}}
\usepackage{langsci-textipa}
\usepackage{langsci-branding}
\usepackage{tabto}

   %%%%%%%%%%%%%%%%%%%%%%%%%%%%%%%%%%%%%%%%%%%%%
%%%%%%%%%% From LSP %%%%%%%%%%%%%%%%%%%%%%%%%
%%%%%%%%%%%%%%%%%%%%%%%%%%%%%%%%%%%%%%%%%%%%%


% \newcommand*{\orcid}{}


\makeatletter
\let\thetitle\@title
\let\theauthor\@author
\makeatother

\newcommand{\togglepaper}[1][0]{
   \bibliography{../localbibliography}
   \papernote{\scriptsize\normalfont
     \theauthor.
     \titleTemp.
     To appear in:
     E. Di Tor \& Herr Rausgeberin (ed.).
     Booktitle in localcommands.tex.
     Berlin: Language Science Press. [preliminary page numbering]
   }
   \pagenumbering{roman}
   \setcounter{chapter}{#1}
   \addtocounter{chapter}{-1}
}

\newbool{bookcompile}
\booltrue{bookcompile}
\newcommand{\bookorchapter}[2]{\ifbool{bookcompile}{#1}{#2}}


%%%%%%%%%%%%%%%%%%%%%%%%%%%%%%%%%%%%%%%%%%%%%
%%%%%%%%%% From ACAL LaTeX template %%%%%%%%%
%%%%%%%%%%%%%%%%%%%%%%%%%%%%%%%%%%%%%%%%%%%%%


\usepackage{tikz}

\newcommand{\circled}[1]{\begin{tikzpicture}[baseline=(word.base)]
\node[draw, rounded corners, text height=8pt, text depth=2pt, inner sep=2pt, outer sep=0pt, use as bounding box] (word) {#1};
\end{tikzpicture}
}


%%%%%%%%%%%%%%%%%%%%%%%%%%%%%%%%%%%%%%%%%%%%%%
%%%%%%%%%%%%%%%%%%%%%%%%%%%%%%%%%%%%%%%%%%%%%%

% Useful Ling Shortcuts


%This makes the \emptyset command be a nicer one
\let\oldemptyset\emptyset
\let\emptyset\varnothing
\newcommand{\nothing}{$\emptyset$}

%Not all of these are explained in the Quick Reference Guide, but they are here bc they are relevant to some of our students. 
\newcommand{\xbar}{\ensuremath{'}\xspace}
\newcommand{\xhead}{\ensuremath{^\circ}\xspace}
\newcommand{\Lb}[1]{$\text{[}_{\text{#1}}$ } %A more convenient left bracket
\newcommand{\Rb}[1]{$\text{]}_{\text{#1}}$ } %A more convenient left bracket
\newcommand{\gap}{\underline{\hspace{1.2em}}}
\newcommand{\vP}{\emph{v}P}
\newcommand{\lilv}{\emph{v}}
\newcommand{\Abar}{A$'$-} %A more convenient A-bar notation
\newcommand{\ph}{$\varphi$\xspace} %A more convenient phi
\newcommand{\pro}{\emph{pro}\xspace}
\newcommand{\subs}[1]{\textsubscript{#1}} %A more convenient subscript
\newcommand{\spells}{$\Longleftrightarrow$} %spellout arrow for morph spellout rules
\newcommand{\tr}[1]{\textit{t}\textsubscript{\textit{#1}}} %easy traces with subscript
\newcommand{\supers}[1]{\textsuperscript{#1}}

%These are borrowed/modified from Andrew Mackenzie's ling-macros package. We have copied them in here (instead of just using the package itself) to avoid a minor clash that was generating an error.


% Abbreviations for glossing, based on Leipzig
\newleipzig{hab}{hab}{habitual}
\newleipzig{rem}{rem}{remote}
\newleipzig{sm}{sm}{subject marker}
\newleipzig{t}{t}{tense}
\newleipzig{aa}{aa}{anti-agreement}
\newleipzig{pron}{pron}{pronoun}
\newleipzig{rec}{rec}{recent}
\newleipzig{om}{om}{object marker}
%\newleipzig{ipfv}{ipfv}{imperfective}
\newleipzig{asp}{asp}{aspect}
\newleipzig{lk}{lk}{linker}
\newleipzig{pcl}{pcl}{particle}
\newleipzig{stat}{stat}{stative}
\newleipzig{ints}{ints}{intensive}
\newleipzig{ascl}{ascl}{assertive subject clitic}
\newleipzig{nascl}{nascl}{non-assertive subject clitic}
\newleipzig{ta}{ta}{tense and/or aspect}
\newleipzig{assoc}{assoc}{associative marker}
\newleipzig{hon}{hon}{honorific}
%\newleipzig{whprt}{wh}{\wh particle}
\newleipzig{sa}{sa}{subject agreement}
\newleipzig{conj}{conj}{conjunction}
%\newleipzig{loc}{loc}{locative}
\newleipzig{expl}{expl}{expletive}
\newleipzig{rcm}{rcm}{reciprocal marker}
\newleipzig{pers}{pers}{persistive}
\newleipzig{fv}{fv}{final vowel}
%\newleipzig{}{}{} %this is just to copy for when I want to add more


%%%%%%%%%%%%%%%%%%%%%%%%%%%%%%%%%%%%%%%%%%%%%%%%%%%%%
%%%%%%%%%%%%%%%%%%%%%%%%%%%%%%%%%%%%%%%%%%%%%%%%%%%%%


%%%%%%%%%%%%%%%%%%%%%%%%%%%%%%%%%%%%%%%%%%%%%
%%%%%%%%%% From Gluckman paper %%%%%%%%%%%%%%
%%%%%%%%%%%%%%%%%%%%%%%%%%%%%%%%%%%%%%%%%%%%%

\newcommand{\pcite}[1]{\citeauthor{#1}'s (\citeyear{#1})}


%%%%%%%%%%%%%%%%%%%%%%%%%%%%%%%%%%%%%%%%%%%%%
%%%%%%%%%% From Myers paper %%%%%%%%%%%%%%
%%%%%%%%%%%%%%%%%%%%%%%%%%%%%%%%%%%%%%%%%%%%%

%% hspace over width of something without showing it
\newcommand*{\hspaceThis}[1]{\settowidth{\LSPTmp}{#1}\hspace*{\LSPTmp}}
% use \hphantom{} for this


%%%%%%%%%%%%%%%%%%%%%%%%%%%%%%%%%%%%%%%%%%%%%
%%%%%%%%%% From Matte paper %%%%%%%%%%%%%%%%%
%%%%%%%%%%%%%%%%%%%%%%%%%%%%%%%%%%%%%%%%%%%%%

%mjkd commented this out in case it's breaking something right now. 
% \newcounter{xlistex}
% \newcommand{\footnoteex}{\stepcounter{xlistex}(\roman{xlistex})}

\newleipzig{sp}{sp}{speaker} %a new leipzig glossing command


%%%%%%%%%%%%%%%%%%%%%%%%%%%%%%%%%%%%%%%%%%%%%
%%%%%%%%%% From GamFranich paper %%%%%%%%%%%%%%
%%%%%%%%%%%%%%%%%%%%%%%%%%%%%%%%%%%%%%%%%%%%%

%MJKD commented these out - both will be provided by the overall LSP book template 

% \bibliography{localbibliography}
% \newcommand{\orcid}[1]{}

%%%%%%%%%%%%%%%%%%%%%%%%%%%%%%%%%%%%%%%%%%%%%
%%%%%%%%%% From Diercks paper %%%%%%%%%%%%%%
%%%%%%%%%%%%%%%%%%%%%%%%%%%%%%%%%%%%%%%%%%%%%

\newcommand{\abar}{$\bar{\text{A}}$\xspace} % this one is different to the \Abar command in Mike's shortcuts doc
\newcommand{\topFeat}{[\textsc{top}]\xspace}


%%%%%%%%%%%%%%%%%%%%%%%%%%%%%%%%%%%%%%%%%%%%%
%%%%%%%%%% From Kinchsular paper %%%%%%%%%%%%%%
%%%%%%%%%%%%%%%%%%%%%%%%%%%%%%%%%%%%%%%%%%%%%

%%%%%%%%%%%%%%%%%%%%%%%%%%%%%%%%%%%%%%%%%%%%%
%%%%%%%%%% From Chen paper %%%%%%%%%%%%%%
%%%%%%%%%%%%%%%%%%%%%%%%%%%%%%%%%%%%%%%%%%%%%

\newcounter{savedex}
\makeatletter
\newcommand*{\saveex}{\setcounter{savedex}{\value{\@xnumctr}}}
\newcommand*{\resumeex}{\setcounter{\@xnumctr}{\value{savedex}}}
\makeatother


   \input{../localhyphenation}
   \boolfalse{bookcompile}
   \togglepaper[23]%%chapternumber
}{}

\begin{document}
\maketitle
\lehead{Matte et al.}
\section{Introduction} \label{sec:matte:1}

Adara (ISO 639-3 [KAD]), also known as Eda, Edra, and Kadara, is an under-documented Benue-Congo language spoken by approximately 300,000 \citep{HonEtAl2018} to 500,000 people \citep{571261} in Kaduna and Niger states in Nigeria, as highlighted in \figref{fig:matte:Figure1}. \figref{fig:matte:Figure2} provides a more detailed map of the Adara-speaking area of Nigeria taken from \citet{HonEtAl2018}, specifying the locations of the various Adara dialects.
\\
\\
\\
\begin{figure}
% 
\includegraphics[scale=0.5]{figures/MatteFigure1.png}
\includegraphics[width=0.8\textwidth]{figures/MatteFig1a.pdf}
\caption{Map of Nigeria showing the Adara-speaking region}
\label{fig:matte:Figure1}
\end{figure}
\begin{figure}
% \includegraphics[scale=0.5]{figures/MatteFigure2.png}
\includegraphics[width=0.8\textwidth]{figures/MatteFigure2.pdf}
\caption{Map of the Adara-speaking area and its dialects (based on \citet{HonEtAl2018})}
\label{fig:matte:Figure2}
\end{figure}

\noindent Despite its large number of speakers, very little research has been done on Adara. This paper presents the first description of question formation in the language. Specifically, we investigate polar and \textit{wh-} questions in Ekhwa Adara\footnote{The data and judgments presented in this paper come exclusively from fieldwork with the third author, a native speaker of Ekhwa Adara, in the context of a Field Methods class at the CUNY Graduate Center in Fall 2022. Data are presented in IPA. The following diacritics are used to mark surface tone: V́ = high, V̀ = low, V = mid; V̂ = falling, V̌ = rising.}, the least researched of Adara’s five dialects. Overall, we have found that polar questions in Ekhwa Adara involve final-vowel lengthening + L\% boundary tones and that \textit{wh-} movement in the language is optional. Where \textit{wh-} movement does occur, the ex-situ \textit{wh-} item obligatorily precedes a corresponding focus marker whose choice is determined by the animacy (human vs. non-human) and grammatical function (argument vs. adjunct) of the \textit{wh-} item. Furthermore, both long-distance and partial \textit{wh-} movement is possible in Ekhwa Adara. Lastly, and perhaps most remarkably, due to its (to the best of our knowledge) unattested status in the grammars of African languages, Ekhwa Adara may allow for multiple \textit{wh-} fronting, but is possibly constrained by an anti-superiority condition.

This article is structured in the following way. In \sectref{sec:matte:2}, we describe the formation of polar questions in Ekhwa Adara. \sectref{sec:matte:3} details our findings on \textit{wh-} questions in Ekhwa Adara, including \textit{wh-} in-situ phenomena, \textit{wh-} movement (including partial and long-distance), embedded question formation, and multiple \textit{wh-} questions. We conclude in \sectref{sec:matte:4} with a summary of our key findings and questions for future research.

\section{Polar questions} \label{sec:matte:2}

Ekhwa Adara’s primary polar question formation strategy is intonational in nature. Matrix polar questions involve final-vowel lengthening and L\% tones, as is characteristic of the languages of the Sudanic belt region \citep{Rialland2007Prosody, Rialland2009, Cahill2012, Cahill2015}. We see that the declarative sentences in \REF{ex:matte:1a}, \REF{ex:matte:2a}, and \REF{ex:matte:new3a} below all contain a short final vowel which is lengthened in their polar interrogative counterparts (\REF{ex:matte:1b}, \REF{ex:matte:2b}, and \REF{ex:matte:new3b}, respectively).\footnote{In response to one reviewer, we are currently unaware of any other prosodic differences between declaratives and interrogatives in Ekhwa Adara.}

\ea \label{ex:matte:1}
\begin{xlist}

\ex \label{ex:matte:1a}
\gll ijâ sù kɹɔ́\\
Ija \Prog{} cry\\
\glt `Ija is crying.'

\ex \label{ex:matte:1b}
\gll ijâ sù kɹɔ̂ː\\
Ija \Prog{} cry.\Q{}\\
\glt `Is Ija crying?'

\end{xlist}
\z

\ea \label{ex:matte:2}
\begin{xlist}

\ex \label{ex:matte:2a}
\gll omúsé ku ɹɔ́ utébur\\
Omuse \Pfv{} buy table\\
\glt `Omuse bought a table.'

\ex \label{ex:matte:2b}
\gll omúsé ku ɹɔ́ utébûːr\\
Omuse \Pfv{} buy table.\Q{}\\
\glt `Did Omuse buy a table?'

\end{xlist}
\z

\ea \label{ex:matte:new3}
\begin{xlist}

\ex \label{ex:matte:new3a}
\gll omúsé ku kɹɔ́ oɹaɪ̀\\
Omuse \Pfv{} cry yesterday\\
\glt `Omuse cried yesterday.'

\ex \label{ex:matte:new3b}
\gll omúsé ku kɹɔ́ oɹaɪ̀ː\\
Omuse \Pfv{} cry yesterday.\Q{}\\
\glt `Did Omuse cry yesterday?'

\end{xlist}
\z

\noindent Furthermore, we see that intonational phrase-final high and mid tones, such as in \REF{ex:matte:1a} \textit{kɹɔ́} `cry' and \REF{ex:matte:2a} \textit{utébur} ‘table’, are both realized with falling pitch movements in their polar interrogative counterparts. In \REF{ex:matte:1b}, the phrase-final syllable \textit{kɹɔ̂ː} exhibits a falling tone and similarly in \REF{ex:matte:2b}, the final syllable of \textit{utébûːr} is also realized with a falling tone. On the other hand, intonational phrase-final L tones such as in \REF{ex:matte:new3a} \textit{oɹaì} ‘yesterday’ are still realized as L in their polar interrogative counterparts, as seen in \REF{ex:matte:new3b}. Crucially, they are not realized with falling pitch movements. We also see that the sentences in \REF{ex:matte:1} make use of the progressive aspect, whereas the sentences in \REF{ex:matte:2} make use of the perfective aspect. The data therefore indicate that final-vowel lengthening and L\% tone insertion are not tense/aspect dependent. Additionally, given that the final vowel quality in each example differs, we conclude that polar questions are formed by true vowel lengthening and not simply the addition of a uniform vowel, which occurs in some languages of the region \citep{Rialland2007, Cahill2015}.

Syntactically, polar questions may be marked by a clause-initial Q particle \textit{kó}, which seems to be borrowed from Hausa \citep{Bawa2023}. While this Q particle is optional in matrix polar questions \REF{ex:matte:3a} and \textit{wh-} in-situ constructions \REF{ex:matte:3c}, it is obligatory in embedded polar questions \REF{ex:matte:3b}.

\ea \label{ex:matte:3}
\begin{xlist}

\ex \label{ex:matte:3a}
\gll (kó) ijâ sù kɹɔ̂ː\\
\Q{} Ija \Prog{} cry.\Q{}\\
\glt `Is Ija crying?'

\ex \label{ex:matte:3c}
\gll (kó) ijâ ku ɹɔ́ ɪncɪ́\\
\Q{} Ija \Pfv{} buy what\\
\glt `What did Ija buy?'

\ex \label{ex:matte:3b}
\gll omúsé ku ɹɔ́:ru *(kó) ijâ ku kɹɔ́\\
Omuse \Pfv{} ask \Q{} Ija \Pfv{} cry\\
\glt `Omuse asked if Ija cried.'

\end{xlist} 
\z

\noindent From \REF{ex:matte:3b}, we also see that in embedded polar questions, where the Q particle obligatorily appears in the embedded clause to mark the scope of the indirect question, there is no final-vowel lengthening or L\% boundary tone. Thus, final-vowel lengthening and insertion of L\% boundary tones represent main clause phenomena in the language, as in Ikpana \citep{KandybowiczKatsuda2023}.

\section{\textit{Wh-} questions} \label{sec:matte:3}

\subsection{\textit{Wh-} in-situ} \label{sec:matte:3.1}

All argument \textit{wh-} items may appear in-situ in root clauses.\footnote{Though as discussed below, it is not clear whether focused subjects are in-situ or ex-situ at this stage of our research.} In \REF{ex:matte:4}, we observe the \textit{wh-} item appearing in subject position \REF{ex: matte:4a}, direct object position (\ref{ex:matte:4b}, \ref{ex:matte:4d}, \ref{ex:matte:4e}), and indirect object position \REF{ex:matte:4c}.\footnote{From (\ref{ex:matte:4c}, \ref{ex:matte:4d}), we observe that in Ekhwa Adara double object constructions, the indirect object immediately precedes the direct object.}

\ea \label{ex:matte:4}
\begin{xlist}

\ex \label{ex: matte:4a}
\gll \textbf{iwé} sú ɹɔ eg͡bé\\
who \Fut{} buy house\\
\glt `Who will buy a house?'

\ex \label{ex:matte:4b}
\gll ijâ ku ɹi \textbf{ɪncɪ́/iwé}\\
Ija \Pfv{} see what/who\\
\glt `What/who did Ija see?'

\ex \label{ex:matte:4c}
\gll ijâ ku dʒe \textbf{iwé} onsɛ̌\\
Ija \Pfv{} give who name\\
\glt `Who did Ija give a name to?'

\ex \label{ex:matte:4d}
\gll ijâ ku dʒe awɛ́-ǹ \textbf{ɪncɪ́}\\
Ija \Pfv{} give child-\Def{} what\\
\glt `What did Ija give the child?'

\ex \label{ex:matte:4e}
\gll ijâ ku ɹɔ́ \textbf{ocɪná} eg͡bé\\
Ija \Pfv{} buy which house\\
\glt `Which house did Ija buy?'

\end{xlist}
\z

As previously discussed in \sectref{sec:matte:2}, the Q particle \textit{kó} may optionally appear in root clause \textit{wh-} in-situ questions. Currently, we only have evidence of its (optional) appearance in object root clause \textit{wh-} in-situ questions such as (\ref{ex:matte:5a}, \ref{ex:matte:5b}). We lack the data to determine whether Q may also appear in non-object root \textit{wh-} in-situ constructions, but we suspect this to be the case.\footnote{Our consultant seems to prefer either ex-situ or focus marked in-situ \textit{wh-} subject questions, at least in root contexts. This may explain why the Q particle is limited to \textit{wh-} object questions in our dataset.} Note that in \REF{ex:matte:5b}, we observe a focused \textit{wh-} subject after \textit{kó}. This, however, should not be considered an instance where \textit{kó} is marking an ex-situ \textit{wh-} subject question, because Q-marking does not occur in ex-situ \textit{wh-} questions (note the absence of the Q particle in the data presented in \sectref{sec:matte:3.2}). Rather, in this sentence, \textit{kó} is marking the in-situ \textit{wh-} object \textit{ɪncɪ́} `what'. Given our consultant’s emphasis on beginning an interrogative sentence with some indication of intent (either through the Q particle \textit{kó} or an initial focused \textit{wh-} phrase), the optional \textit{kó} marker most likely signals that there is (at least) one in-situ interrogative expression following the subject.

\ea \label{ex:matte:5}
\begin{xlist}

\ex \label{ex:matte:5a}
\gll (kó) ijâ ku ɹɔ́ ɪncɪ́\\
\Q{} Ija \Pfv{} buy what\\
\glt `What did Ija buy?'

\ex \label{ex:matte:5b}
\gll (kó) iwé ŋu ku ɹɔ́ ɪncɪ́\\
\Q{} who \Foc{} \Pfv{} buy what\\
\glt `Who bought what?’
\end{xlist}
\z

\noindent All adjunct \textit{wh-} items may also appear in-situ in root clauses. This is demonstrated in \REF{ex:matte:6} for a variety of adjuncts (locative \REF{ex:matte:6a}, temporal \REF{ex:matte:6b}, manner/quantity \REF{ex:matte:6c} and reason \REF{ex:matte:6d}).

\ea \label{ex:matte:6}
\begin{xlist}

\ex \label{ex:matte:6a}
\gll ijâ ku ɹɔ́ eg͡bé ɪmbí\\
Ija \Pfv{} buy house where\\
\glt `Where did Ija buy a house?'

\ex \label{ex:matte:6b}
\gll ijâ ku ɹɔ́ eg͡bé ocɪná atú/apa\\
Ija \Pfv{} buy house which day/time\\
\glt `When (i.e., which day/time) did Ija buy a house?'

\ex \label{ex:matte:6c}
\gll a sù nɛ otúmá-ŋ níní\\
\Tsg{}.\Nom{} \Prog{} do work-\Def{} how\\
\glt `How is s/he doing the work?'\\
`For how much (money) is s/he doing the work?'

\ex \label{ex:matte:6d}
\gll ijâ ku ɹɔ́ eg͡bé-ŋ domɪ́n ɪncɪ́\\
Ija \Pfv{} buy house-\Def{} reason what\\
\glt `Why did Ija buy the house?' 

\end{xlist}
\z

Long-distance \textit{wh-} in-situ is also attested. All \textit{wh-} items may appear in-situ in embedded clausal complements and take wide scope over the matrix clause. We observe this in \REF{ex:matte:7}, where the \textit{wh-} items \textit{iwé} ‘who’ \REF{ex:matte:7a} and \textit{ɪncɪ́} ‘what’ \REF{ex:matte:7b} appear as the subject and object of the embedded clause, respectively. Crucially, these sentences cannot receive indirect question interpretations.\footnote{As discussed in \REF{ex:matte:19} below, this would require relativization.} Instead, the sentences are interpreted as long-distance in-situ interrogatives, with the \textit{wh-} item in the complement clause taking matrix scope.\footnote{The Q particle \textit{kó} is unavailable in this construction to mark the scope of the \textit{wh-} element. In fact, with the Q particle \textit{ko}́ in sentence initial position, the \textit{wh-} questions in \REF{ex:matte:7} become quotative polar questions, expressing ‘Did Omuse say: “Who bought a house?”?’ and ‘Did Omuse say: “What did Ija buy?”?’ respectively.}

\ea \label{ex:matte:7}
\begin{xlist}

\ex \label{ex:matte:7a}
\gll omúsé ga iwé ku ɹɔ́ eg͡bé\\
Omuse say who \Pfv{} buy house\\
\glt `Who did Omuse say bought a house?'\\
NOT: `Omuse said who bought a house.'

\ex \label{ex:matte:7b}
\gll omúsé ga ijâ ku ɹɔ́ ɪncɪ́\\
Omuse say Ija \Pfv{} buy what\\
\glt `What did Omuse say Ija bought?'\\
NOT: `Omuse said what Ija bought.'

\end{xlist}
\z

\noindent Embedded in-situ adjunct \textit{wh-} expressions may only be interpreted as originating in the embedded clause. We demonstrate this in \REF{ex:matte:8}, where the \textit{wh-} items `where,' `when,' `how much,' and `why'  can only be interpreted as being generated in the embedded clause, that is, modifying the verb \textit{ɹɔ́} `buy', as opposed to the matrix verb \textit{ga} `say'.

\ea \label{ex:matte:8} 
\begin{xlist}

\ex \label{ex:matte:8a}
\gll omúsé ga ijâ ku ɹɔ́ eg͡bé ɪmbí\\
Omuse say Ija \Pfv{} buy house where\\
\glt `Where did Omuse say Ija bought a house?'\\
\langscicheckmark `Where-buy'\\
* `Where-say'

\ex \label{ex:matte:8b}
\gll omúsé ga ijâ ku ɹɔ́ eg͡bé ocɪná apa\\
Omuse say Ija \Pfv{} buy house which time\\
\glt `When did Omuse say Ija bought a house?'\\
\langscicheckmark `When-buy'\\
* `When-say'

\ex \label{ex:matte:8c}
\gll omúsé ga ijâ ku ɹɔ́ eg͡bé níní\\
Omuse say Ija \Pfv{} buy house how\\
\glt `For how much did Omuse say Ija bought a house?'\\
\langscicheckmark `For how much-buy'\\
* `For how much-say'

\ex \label{ex:matte:8d}
\gll omúsé ga ijâ ku ɹɔ́ eg͡bé domɪ́n ɪncɪ́\\
Omuse say Ija \Pfv{} buy house reason what\\
\glt `Why did Omuse say Ija bought a house?'\\
\langscicheckmark `Why-buy'\\
* `Why-say'

\end{xlist}
\z

\subsection{\textit{Wh-} movement} \label{sec:matte:3.2}

\subsubsection{Focus particles} \label{sec:matte:3.2.1}

Ekhwa Adara is an optional \textit{wh-} movement language. As we saw in \sectref{sec:matte:3.1}, all \textit{wh-} arguments and adjuncts may remain in-situ with or without the initial Q particle \textit{kó}. 
The interrogative phrase, however, may also be fronted to clause-initial position. This optional \textit{wh-} movement requires a constituent focus particle, which immediately follows the \textit{wh-} item and varies according to the animacy and type of the interrogative expression. When undergoing movement, we see that the human interrogative expression \textit{iwé} `who' must be followed by the particle \textit{ŋu} \REF{ex:matte:9a}\footnote{Unlike with long-distance \textit{wh-} movement, there is no direct evidence in cases of monoclausal \textit{wh-} questions that the subject is moved at all, since \textit{wh-} items may be accompanied by a focus marker even when left in-situ, as shown in \REF{ex:matte:15}. There also seems to be a preference for the focus-marked construction when the interrogative phrase is the subject, while the same is not true for other arguments. Our point, in presenting \REF{ex:matte:9a}, is merely that, if moved, the \textit{wh-} subject should behave like non-subject \textit{wh-} items and require the focus particle.}, the non-human interrogative expression \textit{ɪncɪ́} `what' must be followed by the particle \textit{mo} \REF{ex:matte:9b}, and finally, the D-linked interrogative \textit{ocɪná} \REF{ex:matte:9g} `which' and interrogative adjuncts such as \textit{ɪmbí} `where' \REF{ex:matte:9c}, \textit{ocɪná apa} `when' \REF{ex:matte:9d}, and \textit{níní} `how/for how much' \REF{ex:matte:9e} must be followed by the particle \textit{ku}. The only exception is the \textit{wh-} adjunct \textit{domɪ́n ɪncɪ́} `why' (literally, `what reason') which may either precede \textit{ku}, like all other \textit{wh-} adjuncts, or \textit{mo}, like \textit{ɪncɪ́} `what' \REF{ex:matte:9f}.

\ea \label{ex:matte:9}
\begin{xlist}

\ex \label{ex:matte:9a}
\gll iwé *(ŋu) {\longrule} ku ɹɔ́ eg͡bé\\
who \Foc{} {} \Pfv{} buy house\\
\glt `Who bought a house?'

\ex \label{ex:matte:9b}
\gll ɪncɪ́ *(mo) ijâ ku ɹɔ́ {\longrule}\\
what \Foc{} Ija \Pfv{} buy {}\\
\glt `What did Ija buy?'

\ex \label{ex:matte:9g}
\gll ocɪná egbé *(ku) ijâ ku ɹɔ́ {\longrule}\\
which house \Foc{} Ija \Pfv{} buy {}\\
\glt `Which house did Ija buy?'

\ex \label{ex:matte:9c}
\gll ɪmbí *(ku) ijâ ku ɹɔ́ egbé {\longrule}\\
what \Foc{} Ija \Pfv{} buy house {}\\
\glt `Where did Ija buy a house?'

\ex \label{ex:matte:9d}
\gll ocɪná apa *(ku) ijâ ku ɹɔ́ egbé {\longrule}\\
which time \Foc{} Ija \Pfv{} buy house {}\\
\glt `When did Ija buy a house?'

\ex \label{ex:matte:9e}
\gll níní *(ku) a sù nɛ otúmá-ŋ {\longrule}\\
how \Foc{} \Tsg{}.\Nom{} \Prog{} do work-\Def{}\\
\glt `How/for how much is s/he doing the work?'

\ex \label{ex:matte:9f}
\gll domɪ́n ɪncɪ́ *(ku/mo) ijâ ku ɹɔ́ eg͡bé {\longrule}\\
reason what \Foc{} Ija \Pfv{} buy house\\
\glt `Why did Ija buy a house?'


\end{xlist}
\z

The ungrammaticality of the examples in \REF{ex:matte:10}, in which the focus particles that follow the \textit{wh-} phrase are switched, demonstrate that these markers are compatible only with specific \textit{wh-} items – human \textit{iwé} `who' must be paired with \textit{ŋu}, rather than \textit{mo} or \textit{ku}; non-human \textit{ɪncɪ́} `what' must be paired with \textit{mo}, rather than \textit{ŋu} or \textit{ku}; lastly, the D-linked interrogative \textit{ocɪná} `which' and \textit{wh-} adjuncts, with the exception of \textit{domɪ́n ɪncɪ́} `why', must be paired with \textit{ku}, rather than \textit{ŋu} or \textit{mo}. 

\ea \label{ex:matte:10}
\begin{xlist}

\ex \label{ex:matte:10a} 
\gll * iwé mo/ku {\longrule} ku ɹɔ́ eg͡bé\\
{} who \Foc{} {} \Pfv{} buy house\\
\glt Intended: `Who bought a house?'

\ex \label{ex:matte:10b}
\gll * ɪncɪ́ ŋu/ku ijâ ku ɹɔ́ {\longrule}\\
{} what \Foc{} Ija \Pfv{} buy {}\\
\glt Intended: `What did Ija buy?'

\ex \label{ex:matte:10g}
\gll * ocɪná egbé ŋu/mo ijâ ku ɹɔ́ {\longrule}\\
{} which house \Foc{} Ija \Pfv{} buy {}\\
\glt Intended: `Which house did Ija buy?'

\ex \label{ex:matte:10c}
\gll * ɪmbí ŋu/mo ijâ ku ɹɔ́ egbé {\longrule}\\
{} what \Foc{} Ija \Pfv{} buy house {}\\
\glt Intended: `Where did Ija buy a house?'

\ex \label{ex:matte:10d}
\gll * ocɪná apa ŋu/mo ijâ ku ɹɔ́ egbé {\longrule}\\
{} which time \Foc{} Ija \Pfv{} buy house {}\\
\glt Intended: `When did Ija buy a house?'

\ex \label{ex:matte:10e}
\gll * níní ŋu/mo a sù nɛ otúmá-ŋ {\longrule}\\
{} how \Foc{} \Tsg{}.\Nom{} \Prog{} do work-\Def{}\\
\glt Intended: `How/for how much is s/he doing the work?'

\ex \label{ex:matte:10f}
\gll * domɪ́n ɪncɪ́ ŋu ijâ ku ɹɔ́ eg͡bé {\longrule}\\
{} reason what \Foc{} Ija \Pfv{} buy house\\
\glt Intended: `Why did Ija buy a house?'

\end{xlist}
\z

Furthermore, the examples in \REF{ex:matte:11} show that the choice of \textit{ŋu} or \textit{mo} is truly dependent on the interrogative item itself and not on its syntactic relationship to the verb (e.g., case). As we can see, ex-situ \textit{iwé} still requires \textit{ŋu} when thematically linked to non-subject positions like direct objects \REF{ex:matte:11a} and possessors of objects \REF{ex:matte:11b}. Likewise, ex-situ \textit{ɪncɪ́} selects the marker \textit{mo} even when linked to the subject position \REF{ex:matte:11c}. 

\ea \label{ex:matte:11}
\begin{xlist}

\ex \label{ex:matte:11a}
\gll iwé ŋu ijâ	ku dá {\longrule}\\
who \Foc{} Ija \Pfv{} touch\\
\glt `Who has Ija touched?'

\ex \label{ex:matte:11b}
\gll iwé ŋu	ijâ	ku ɹɔ́ eg͡bé-ŋ {\longrule}\\
who	\Foc{} Ija \Pfv{} buy house-\Def{}\\
\glt `Who did Ija buy a house from?'

\ex \label{ex:matte:11c}
\gll ɪncɪ́ mo {\longrule} ku dá ijâ\\
what \Foc{} {} \Pfv{} touch Ija\\
\glt `What has touched Ija?'

\end{xlist}
\z

\noindent In addition to appearing in ex-situ \textit{wh-} questions, these constituent focus particles may occur in answer fragments. Here, we see that \REF{ex:matte:12b} is a grammatical response to \REF{ex:matte:12a}.


\ea \label{ex:matte:12}
\begin{xlist}

\ex \label{ex:matte:12a}
\gll iwé ŋu mɪ́n\\
who \Foc{} here\\
\glt `Who is this?'

\ex \label{ex:matte:12b}
\gll omúsé ŋu\\ 
Omuse \Foc{}\\
\glt `(This is) Omuse.'

\end{xlist}
\z

Both \textit{ŋu} and \textit{ku}\footnote{The constituent focus particles \textit{ŋu} and \textit{ku} also occur in allomorphic variation with \textit{ŋo} and \textit{ko}, respectively. Further research is necessary to see whether this variation is phonologically conditioned. Note that the constituent focus particle variant \textit{ko}, which bears a mid tone, is tonally (not segmentally) contrastive with the high tone-bearing initial interrogative particle \textit{kó}.} also appear as focus particles in declarative sentences with focus fronting. Like in ex-situ interrogative constructions, focus markers follow left peripheral foci and the realization of the surfacing focus particle depends on the constituent it focus-marks. Consider the facts in \REF{ex:matte:13} and \REF{ex:matte:14}, which illustrate.\footnote{We initially glossed (\ref{ex:matte:13}-\ref{ex:matte:14}) with English \textit{it}-cleft translations (e.g., \REF{ex:matte:13a} was glossed as ‘It is \textsc{ija} who bought a house’). Thanks to one of our reviewers, however, we realized that these translations might not be semantically appropriate given that \textit{it}-cleft sentences in English have exhaustive properties which the sentences in (\ref{ex:matte:13}-\ref{ex:matte:14}) do not possess. For example, our consultant notes that in the case where Omuse will teach both Ija and Awen a song, \REF{ex:matte:13b} can only be used if the hearer is unaware of Awen's existence. We assume this does not constitute an exhaustive interpretation since Awen's existence is not part of the common ground shared by both speakers. Thus, Awen is not in the domain of individuals whose "learning of the song from Omuse" we need to negate. An exhaustive interpretation of \REF{ex:matte:13b} is only available when both Ija and Awen are relevant to the context (i.e., in the common ground) but only Ija will be taught the song (possibly among other students who, similar to Awen in the previous case, are not part of the common ground and thus, the alternative propositions involving them need not be negated).}

\ea \label{ex:matte:13}
\begin{xlist}

\ex \label{ex:matte:13a}
\gll ijâ ŋu/*ku ɹɔ́ eg͡bé\\
Ija \Foc{} buy house\\
\glt `\textsc{ija} bought a house.'

\ex \label{ex:matte:13b}
\gll ijâ ŋu/*ku	omúsé sú mani {\longrule} ɔnʃó\\
Ija \Foc{} Omuse \Fut{} teach {} song\\
\glt `Omuse will teach a song (to) \textsc{ija}.'

\ex \label{ex:matte:13c}
\gll ená ŋu/??ku omúsé sú mani {\longrule} ɔnʃó\\
cow	\Foc{} Omuse \Fut{} teach	{} song\\
\glt `Omuse will teach a song to a \textsc{cow}.'

\end{xlist}
\z

\ea \label{ex:matte:14}
\begin{xlist}

\ex \label{ex:matte:14a}
\gll ɔnʃó ku/*ŋu omúsé sú mani-ŋ ijâ {\longrule}\\
song \Foc{} Omuse \Fut{} teach-\Tsg{}.\Acc{} Ija {}\\
\glt `Omuse will teach Ija a \textsc{song}.'

\ex \label{ex:matte:14b}
\gll eg͡bé ku/*ŋu/*mo ijâ ku ɹɔ́ {\longrule}\\
house \Foc{} Ija \Pfv{} buy {}\\
\glt `Ija bought a \textsc{house}.'

\end{xlist}
\z

\noindent The data in \REF{ex:matte:13} and \REF{ex:matte:14} suggest that the choice of focus marker in declarative sentences is related to animacy – focused animate DPs appear with \textit{ŋu} \REF{ex:matte:13}, while focused inanimates are accompanied by \textit{ku} \REF{ex:matte:14}. As before with ex-situ \textit{wh-} questions, the choice of either \textit{ŋu} or \textit{ku} in declarative sentences does not seem to depend on the relationship of the moved phrase to other elements of the sentence (e.g., case, thematic role, etc.).

The focus markers introduced in this section are not limited to contexts in which an interrogative undergoes movement. \textit{Wh-} items that remain in-situ may optionally precede these same particles \REF{ex:matte:15}. 

\ea \label{ex:matte:15}
\begin{xlist}

\ex \label{ex:matte:15a}
\gll ijâ ku dʒe iwé (ŋu) onsɛ̌\\
Ija \Pfv{} give who \Foc{} name\\
\glt `Who did Ija give a name to?'
 
\ex \label{ex:matte:15b}
\gll ijâ ku ɹi ɪncɪ́ (mo)\\
Ija \Pfv{} see what \Foc{}\\
\glt `What did Ija see?'

\ex \label{ex:matte:15c}
\gll ijâ ku ɹɔ́ eg͡bé ɪmbí (ku)\\
Ija \Pfv{} buy house where \Foc{}\\
\glt `Where did Ija buy a house?'

\ex \label{ex:matte:15d}
\gll ijâ ku ɹɔ́ eg͡bé ocɪná apa (ku)\\
Ija \Pfv{} buy house which time \Foc{}\\
\glt `When did Ija buy a house?'

\ex \label{ex:matte:15e}
\gll a sù nɛ otúmá-ŋ níní (ku)\\
\Tsg{}.\Nom{} \Prog{} do work-\Def{} how \Foc{}\\
\glt `How/for how much is s/he doing the work?'

\ex \label{ex:matte:15f}
\gll ijâ ku ɹɔ́ eg͡bé-ŋ domɪ́n ɪncɪ́ (ku/mo)\\
Ija \Pfv{} buy house-\Def{} reason what \Foc{}\\
\glt `Why did Ija buy the house?'

\end{xlist}
\z

A list of Ekhwa Adara \textit{wh-} items and the focus particles they are compatible with is presented in \tabref{tab:matte:1} below.

\begin{table}
\caption{Ekhwa Adara \textit{wh-} items and their accompanying focus particles}
\label{tab:matte:1}

\begin{tabular}{ll}
\lsptoprule
    \textit{Wh-} item & Accompanying focus particle\\
    \midrule  
    \textit{iwé} `who' & \textit{ŋu}\\
    \textit{ɪncɪ́} `what' & \textit{mo}\\
    \textit{ɪmbí} `where' & \textit{ku}\\
    \textit{ocɪná} `which' & \textit{ku}\\
    \textit{ocɪná atú/apa} `when' & \textit{ku}\\
    \textit{níní} `how' & \textit{ku}\\
    \textit{domɪ́n ɪncɪ́} `why' & \textit{ku/mo}\\
\lspbottomrule
\end{tabular}
\end{table}

\subsubsection{Long-distance movement} \label{sec:matte:3.2.2}

Long-distance \textit{wh-} movement in Ekhwa Adara exhibits the same dependency between the moved phrase and its accompanying focus marker. Additionally, however, we see an asymmetry between subjects and non-subjects. The subject of an embedded clause, when fronted to sentence-initial position, must be resumed by an agreeing pronominal \REF{ex:matte:16a}, while resumption of displaced embedded objects is ungrammatical \REF{ex:matte:16b}.\footnote{We did not observe resumption in the indirect object or adjunct positions in long-distance \textit{wh-} movement. Furthermore, to the best of our knowledge, the language does not have any null resumptive forms. That is, it only has overt pronominals.} This is typical of languages in the region, such as Nupe \citep{Kandybowicz2004, Kandybowicz2008} and Igbo \citep{AmaechiGeorgi2019}. 

\ea \label{ex:matte:16}
\begin{xlist}

\ex \label{ex:matte:16a}
\gll iwé ŋu omúsé ga *(a) ku ɹɔ́ eg͡bé\\
who \Foc{} Omuse say \Tsg{}.\Nom{} \Pfv{} buy house\\
\glt `Who did Omuse say bought a house?'

\ex \label{ex:matte:16b}
\gll ɪncɪ́ mo omúsé ga ijâ ku ɹɔ́ (*ŋ)\\
what \Foc{} Omuse say Ija \Pfv{} buy \Tsg{}.\Acc{}\\
\glt `What did Omuse say Ija bought?'

\end{xlist}
\z


\noindent Examples \REF{ex:matte:16c}-\REF{ex:matte:16e} demonstrate that movement of a \textit{wh-} adjunct in biclausal questions does not require pronominal resumption or any other indication of the position in which the \textit{wh-} adjunct was generated. Accordingly, the sentences may receive either matrix or embedded interpretations, questioning the saying event or the buying event, respectively. This is in contrast to examples containing in-situ \textit{wh-} adjuncts, as well as in-situ \textit{wh-} arguments. As we saw earlier in \REF{ex:matte:8}, long-distance (i.e., embedded) \textit{wh-} in-situ with adjuncts forces embedded interpretations, meaning that the \textit{wh-} phrase can only be interpreted as modifying the lower clause verb, whereas in \REF{ex:matte:7}, long-distance \textit{wh-} in-situ with arguments forces matrix interpretations, meaning that such sentences must be interpreted as long-distance in-situ \textit{wh-} interrogatives as opposed to indirect questions.

\ea \label{ex:matte:split1718}

\begin{xlist}

    \ex \label{ex:matte:16c}
    \gll ɪmbí ku omúsé ga ijâ ku ɹɔ́	eg͡bé\\
    where \Foc{} Omuse say Ija \Pfv{} buy house\\
    \glt `Where did Omuse say Ija bought a house?'\\
    \langscicheckmark `Where-buy'\\
    \langscicheckmark `Where-say'

    \ex \label{ex:matte:16d}
    \gll ocɪná apa ku omúsé ga ijâ ku ɹɔ́ eg͡bé\\
    which time \Foc{} Omuse say Ija \Pfv{} buy house\\
    \glt `When did Omuse say Ija bought a house?'\\
    \langscicheckmark `When-buy'\\
    \langscicheckmark `When-say'
    
    \ex \label{ex:matte:16e}
    \gll domɪ́n ɪncɪ́ ku/mo omúsé ga ijâ ku ɹɔ́ eg͡bé\\
    reason what \Foc{} Omuse say Ija \Pfv{} buy house\\
    \glt `Why did Omuse say Ija bought a house?'\\
    \langscicheckmark `Why-buy'\\
    \langscicheckmark `Why-say'
    
\end{xlist}
\z

\subsubsection{Partial movement} \label{sec:matte:3.2.3}

Ekhwa Adara also allows partial \textit{wh-} movement, that is, movement of an embedded interrogative to the initial position of a lower clause. When partially moved, \textit{wh-} is accompanied by its corresponding focus particle (as discussed in \sectref{sec:matte:3.2.1}). Unlike in long-distance \textit{wh-} questions, partially moved subjects in Ekhwa Adara do not seem to require pronominal resumption \REF{ex:matte:17a}, erasing the subject–non-subject asymmetry that was exhibited in (\ref{ex:matte:16a}, \ref{ex:matte:16b}) above.

\ea \label{ex:matte:17}
\begin{xlist}

\ex \label{ex:matte:17a}
\gll omúsé ga iwé ŋu {\longrule} ku ɹɔ́ eg͡bé\\
Omuse say who \Foc{} {} \Pfv{} buy house\\
\glt `Who did Omuse say bought a house?'

\ex \label{ex:matte:17b}
\gll omúsé ga ɪncɪ́ mo ijâ ku ɹɔ́ {\longrule}\\
Omuse say what \Foc{} Ija \Pfv{} buy {}\\
\glt `What did Omuse say Ija bought?'

\ex \label{ex:matte:17c}
\gll omúsé ga ɪmbí ku ijâ ku ɹɔ́ eg͡bé {\longrule}\\
Omuse say where \Foc{} Ija \Pfv{} buy house {}\\
\glt `Where did Omuse say Ija bought a house?'\\
\langscicheckmark `Where-buy'\\
* `Where-say'

\ex \label{ex:matte:17d}
\gll omúsé ga ocɪná apa ku ijâ ku ɹɔ́ eg͡bé {\longrule}\\
Omuse say which time \Foc{} Ija \Pfv{} buy house\ {}\\
\glt `When did Omuse say Ija bought a house?'\\
\langscicheckmark `When-buy'\\
* `When-say'

\ex \label{ex:matte:17e}
\gll omúsé ga níní ku ijâ ku ɹɔ́ eg͡bé {\longrule}\\
Omuse say how \Foc{} Ija \Pfv{} buy house {}\\
\glt `For how much did Omuse say Ija bought a house?'\\
\langscicheckmark `How much-buy'\\
* `How much-say'

\ex \label{ex:matte:17f}
\gll omúsé ga domɪ́n ɪncɪ́ mo ijâ ku ɹɔ́ eg͡bé {\longrule}\\
Omuse say reason what \Foc{} Ija \Pfv{} buy house {}\\
\glt `Why did Omuse say Ija bought a house?'\\
\langscicheckmark `Why-buy'\\
* `Why-say'

\end{xlist}
\z

\noindent As with embedded in-situ questions, partial movement of the \textit{wh-} adjuncts in \REF{ex:matte:17c}–\REF{ex:matte:17f} can only be interpreted as originating in the immediately-containing clause (at least in biclausal contexts). Questions formed in this way do not require long-distance movement to sentence-initial position to mark scope, nor do they require (or allow) the clause-initial Q marker \textit{kó}. We can thus classify Ekhwa Adara as a Simple/Naked partial \textit{wh-} movement language, referencing \cite{Fanselow2006}'s typology. In this way, partial \textit{wh-} movement in Ekhwa Adara resembles partial \textit{wh-} fronting in other West African languages such as Krachi \citep{TorrenceKandybowicz2013, TorrenceKandybowicz2015, KandybowiczTorrence2015}; Bono and Wasa \citep{Kandybowicz2017, Kandybowicz2020, KandybowiczTorrence2015}; Avatime \citep{DevlinTorrence2021}; and Ikpana \citep{KandybowiczKatsuda2021, KandybowiczKatsuda2023}. 

Interestingly, there is seemingly a limit on how far a \textit{wh-} item can move in partial movement structures. In \REF{ex:matte:new20}, where there are two levels of embedding, we see that \textit{wh-} items can only either partially move within the originating clause \REF{ex:matte:new20a} without the resumptive pronoun or move long-distance to the matrix operator position \REF{ex:matte:new20b} with the resumptive pronoun. They cannot move to any intermediate position. That is, a \textit{wh-} item originating in the lowest embedded clause cannot move to the left periphery of an intermediate clause, with or without a resumptive pronoun \REF{ex:matte:new20c}.

\ea \label{ex:matte:new20}

\begin{xlist}

\ex \label{ex:matte:new20a}
\gll awɛ̂n ga (aɪ) ijâ ga-ŋ (aɪ)	iwé ŋu {\longrule} ku ɹɔ́ eg͡be\\
Awen say \Comp{} Ija say-\Tsg{}.\Acc{} \Comp{} who \Foc{} {} \Pfv{} buy house\\
\glt `Who did Awen say that Ija said bought a house?'\\
Lit: `Who did Awen say that Ija told her (Awen) bought a house?'


\ex \label{ex:matte:new20b}
\gll iwé ŋu awɛ̂n ga (aɪ) ijâ ga-ŋ (aɪ) a ku ɹɔ́ eg͡be\\
who \Foc{} Awen say \Comp{} Ija say-\Tsg{}.\Acc{} \Comp{} \Tsg{}.\Nom{} \Pfv{} buy house\\
\glt `Who did Awen say that Ija said bought a house?'\\
Lit: `Who did Awen say that Ija told her (Awen) bought a house?'

\ex \label{ex:matte:new20c}
\gll *awɛ̂n ga (aɪ) iwé ŋu ijâ ga-ŋ (aɪ) (a) ku ɹɔ́ eg͡be\\
Awen say \Comp{} who \Foc{} Ija say-\Tsg{}.\Acc{} \Comp{} \Tsg{}.\Nom{} \Pfv{} buy house\\
\glt Intended: `Who did Awen say that Ija said bought a house?'

\end{xlist}
\z

\subsection{Embedded questions} \label{sec:matte:3.3}

\textit{Wh-} movement to the left periphery of the lower clause also occurs in the formation of embedded questions under the matrix verb `ask' \REF{ex:matte:18}.\footnote{The 
    lower clauses in \REF{ex:matte:18} are indeed embedded questions under the matrix verb `ask', rather than, as per a reviewer’s concern, quotations. In \REF{ex:matte:i}, the embedded first-person object refers to the speaker, as expected of an embedded question interpretation, rather than the matrix subject Omuse, which is what we would expect of a quotative interpretation. 

    \ea \label{ex:matte:i}
        \gll omúsé\textsubscript{i}	ku	ɹɔ́:ru iwé	ŋu	{\longrule} ri-i\textsubscript{\Sp{}/*i}\\
    	Omuse \Pfv{} ask who \Foc{}	{} see-\Fsg{}.\Acc{}\\
        \glt `Omuse\textsubscript{i} asked who saw me\textsubscript{\Sp{}/*i}'
    
    \z
} This construction is identical to the partial \textit{wh-} movement construction described above – both subjects and objects are moved without use of a resumptive pronoun and the clause-initial Q particle is absent. The scope of the \textit{wh-} item in this case is limited to the lower clause due (presumably) to the semantic properties of `ask.'

\ea \label{ex:matte:18}
\begin{xlist}

\ex \label{ex:matte:18a}
\gll omúsé ku ɹɔ́:ru iwé ŋu {\longrule} ku ɹɔ́ eg͡bé\\
Omuse \Pfv{} ask who \Foc{} {} \Pfv{} buy house\\
\glt `Omuse asked who bought a house.'

\ex \label{ex:matte:18b}
\gll omúsé ku ɹɔ́:ru ɪncɪ́ mo ijâ ku ɹɔ́ {\longrule}\\
Omuse \Pfv{} ask what \Foc{} Ija \Pfv{} buy {}\\
\glt `Omuse asked what Ija bought.'

\end{xlist}
\z

Indirect questions with `ask' do not, however, require \textit{wh-} movement in the embedded clause. \REF{ex:matte:new22} shows an example of an embedded \textit{wh-} in-situ question under `ask'. 

\ea \label{ex:matte:new22}

\gll omúsé ku ɹɔ́:ru-i	ijâ	ku	rɔ́ ɪncɪ́ mo\\
Omuse \Pfv{} ask-\Fsg{} Ija \Pfv{} buy what \Foc{}\\
\glt `Omuse asked me what Ija bought.'\footnote{Our consultant notes that in \REF{ex:matte:new22}, without the matrix first-person object \textit{i}, the speaker would expect an answer (except in certain contexts), which defeats the purpose of the reported speech we are aiming to document.}

\z

Indirect question complements of other embedding verbs are formed via relativization, generating what can be described as “concealed questions”. The sentence in \REF{ex:matte:19a} is nearly identical to that in \REF{ex:matte:18a}, apart from the matrix verb and aspect marking\footnote{In many cases, we found a preference for certain sentences to be produced without overt aspect marking, even if the corresponding sentence with a marker was judged to be grammatical. Absence of the aspectual element in \REF{ex:matte:19a}, therefore, would not be expected to negatively affect its acceptability. If anything, we might expect it to improve it.}; however, only \REF{ex:matte:18a} is grammatical. Local fronting of the \textit{wh-} item within the lower clause does not achieve the intended reading in \REF{ex:matte:19a} – i.e., it is not interpreted as an indirect question.\footnote{Currently, we are unsure whether \REF{ex:matte:19a} can be parsed as an instance of partial \textit{wh-} movement, with the moved phrase taking scope over the matrix clause. Given that `know' can embed a clause (see \REF{ex:matte:ii} in footnote 15 below), if the partial \textit{wh-} parse is not available, this may suggest that `know' predicates in Ekhwa Adara can form factive islands.} In order to attain the intended reading, it is necessary to use a relative clause construction. In the case of an embedded `who' question, the relative clause is headed by \textit{anú-ǹ} `the person' \REF{ex:matte:19b}, replacing the \textit{wh-} phrase \textit{iwé ŋu} `who'. In the case of an indirect `what' question, the relative clause is headed by \textit{ɔŋgǎ-ŋ̀} ‘the thing’ \REF{ex:matte:19c} replacing \textit{ɪncɪ́ mo} `what'.\footnote{Although
    the verb `know' is followed by a noun phrase in every example in \REF{ex:matte:19}, we do not analyze \textit{k͡péɹi} `know' as an acquaintance verb similar to French \textit{connaitre} because it can be followed by a clause, as seen in \REF{ex:matte:ii}.

    \ea \label{ex:matte:ii}
        \gll omúsé	k͡péɹi	aɪ	ijâ	ɹɔ́	egbé\\ 
        Omuse know \Comp{} Ija buy house\\
        \glt `Omuse knows that Ija bought a house.'
    \z
}
As indicated below, these structures translate more literally as `the person/thing that...'. Subject and object relative clauses of this type are both generally grammatical in Ekhwa Adara, using this exact construction.


\ea \label{ex:matte:19}
\begin{xlist}

\ex \label{ex:matte:19a}
\gll * omúsé k͡péɹi iwé ŋu ku ɹɔ́ eg͡bé\\
{} Omuse know who \Foc{} \Pfv{} buy house\\
\glt Intended: `Omuse knows who bought a house.'

\ex \label{ex:matte:19b}
\gll omúsé k͡péɹi anú-ǹ da ku ɹɔ́ eg͡bé\\
Omuse know person-\Def{} \Rel{} \Pfv{} buy house\\
\glt `Omuse knows who bought a house.'\\
Literally: `Omuse knows the person who bought a house.'

\ex \label{ex:matte:19c}
\gll omúsé k͡péɹi ɔŋgǎ-ŋ̀ da ijâ ku ɹɔ́\\
Omuse know thing-\Def{} \Rel{} Ija \Pfv{} buy\\
\glt `Omuse knows what Ija bought.'\\
Literally: `Omuse knows the thing that Ija bought.'

\end{xlist}
\z

\noindent This limitation of true embedded questions to `ask' complement clauses connects Ekhwa Adara to Krachi \citep{TorrenceKandybowicz2015}, where identical facts obtain.

\subsection{Multiple \textit{wh-} questions} \label{sec:matte:3.4}

At present, we have observed three potential strategies for the formation of multiple \textit{wh-} questions in Ekhwa Adara: (i) all \textit{wh-} items remain in-situ; (ii) one \textit{wh-} item is fronted to clause-initial position, while the other(s)\footnote{We do not have any data pertaining to questions involving more than two \textit{wh-} elements, though we suspect that such questions would also use this strategy.} remain(s) in-situ; and (iii) all \textit{wh-} items are fronted (i.e., multiple \textit{wh-} fronting). As expected, focus particles appearing in each of these three construction types must be the ones selected by the accompanying \textit{wh-} item. We are certain that strategies (i) and (ii) exist in the grammar. We are less certain about the reality of strategy (iii) given unstable judgments observed over a one-year period, which we discuss below. 

Beginning with strategy (i), example \REF{ex:matte:20} demonstrates that multiple \textit{wh-} questions may be formed simply by leaving all interrogative expressions in-situ.

\ea \label{ex:matte:20}
\gll (*kó) omúsé sú mani iwé (ŋu) ɪncɪ́ (mo)\\
\Q{} Omuse \Fut{} teach who \Foc{} what \Foc{}\\
\glt `What will Omuse teach to who?'
\z

\noindent Note that multiple \textit{wh-} questions formed in this way may not be marked by the clause-initial Q particle \textit{kó}\footnote{A sequence of two clause-initial occurrences of \textit{kó} (i.e., one for each in-situ \textit{wh-} item) is also prohibited in structures like \REF{ex:matte:20}.}, despite the fact that the Q marker may optionally appear when a single \textit{wh-} item is left in-situ \REF{ex:matte:5}. Although focus markers may optionally accompany each in-situ interrogative expression, there is a preference for realizing the focus particle after the first \textit{wh-} item, but not the second. 

Turning next to strategy (ii), example \REF{ex:matte:21} demonstrates that multiple \textit{wh-} questions formed by moving a single interrogative item may optionally be marked by the Q particle \textit{kó}, unlike certain other types of \textit{wh-} movement constructions in Ekhwa Adara discussed in this paper (e.g., long-distance \textit{wh-} movement, partial \textit{wh-} movement, and possibly multiple \textit{wh-} fronting, as described below). In example \REF{ex:matte:21}, we also observe that the hierarchically inferior \textit{wh-} item (the object `what') is fronted, crossing over the higher in-situ interrogative (the subject `who'), an apparent violation of the superiority condition \citep{KunoRobinson1972, Chomsky1973, Chomsky1977}. 

\ea \label{ex:matte:21}
\gll (kó) ɪncɪ́ mo iwé ku ɹɔ́	{\longrule}\\
\Q{} what \Foc{} who \Pfv{} buy {}\\ 
\glt `Who bought what?'
\z

\noindent Similar behavior in the distribution of moved interrogative phrases in multiple \textit{wh-} fronting constructions (see below) also suggests an absence of superiority effects in the language. In this way, Ekhwa Adara is similar to a number of other West African languages that have been shown to lack superiority effects outright \citep{Saah1994, Adesola2005, Adesola2006, TorrenceKandybowicz2015, KandybowiczKatsuda2023, SchurrTapia2023}.

The sentences in \REF{ex:matte:22} below demonstrate strategy (iii) – fronting of multiple \textit{wh-} phrases.\footnote{Since focused in-situ subjects are indistinguishable from their moved counterparts, we avoid presenting multiple \textit{wh-} movement structures in which one interrogative is the subject. We attempted to elicit multiple \textit{wh-} fronting structures in which one of the moved elements was an adjunct, but the types of marginal judgements that we received for these require further investigation.} When more than one \textit{wh-} item is moved to initial position\footnote{We have only collected monoclausal multiple \textit{wh-} questions so far. `Initial' here, therefore, means both clause- and sentence-initial, but we have not yet determined that this must be the case.}, the items are pronounced in an order inverse to the hierarchical position of their extraction sites. 

\ea \label{ex:matte:22}
\begin{xlist}

\ex \label{ex:matte:22a}
\gll omúsé sú mani ijâ ɔnʃó\\
Omuse \Fut{} teach Ija song\\
\glt `Omuse will teach Ija a song.'

\ex \label{ex:matte:22b}
\gll ɪncɪ́\textsubscript{i} mo iwé\textsubscript{j} ŋu omúsé sú mani-*(ŋ)\textsubscript{j} t\textsubscript{i}\\
what \Foc{} who \Foc{} Omuse \Fut{} teach-\Tsg{}.\Acc{} {}\\
\glt `What will Omuse teach to who?'

\ex \label{ex:matte:22c}
\gll * iwé\textsubscript{j} ŋu ɪncɪ́\textsubscript{i} mo omúsé sú mani-*(ŋ)\textsubscript{j} t\textsubscript{i}\\
{} who \Foc{} what \Foc{} Omuse \Fut{} teach-\Tsg{}.\Acc{} {}\\
\glt Intended: `What will Omuse teach to who?'

\end{xlist}
\z

\noindent The canonical ordering of the two lower arguments in a double-object construction places the indirect object in a position immediately preceding the direct object \REF{ex:matte:22a}. When both of these arguments are fronted, as in \REF{ex:matte:22b}, however, the direct object interrogative (`what') must precede the indirect object (`who'). The reverse order, which would both respect the superiority condition and remain closer in form to the canonical Ekhwa Adara word order, was consistently judged ungrammatical \REF{ex:matte:22c}. Assuming the accuracy of these judgment contrasts, we therefore observe a kind of anti-superiority effect in Ekhwa Adara multiple \textit{wh-} movement constructions – the initial fronted \textit{wh-} item must be hierarchically inferior to the other(s). The acceptability of sentences like \REF{ex:matte:22b} might also suggest a recursive FocP domain in the Ekhwa Adara left periphery (contra \citet{Rizzi1997}), unlike what we see in other African languages (e.g., \citealt{aboh2004b}). However, the availability of focus-marking on in-situ \textit{wh-} elements in the language (see \REF{ex:matte:15}) likely supports an alternative analysis in which the peripheral focus particles that surface in multiple \textit{wh-} fronting structures are more accurately “term focus” markers. 

Further research is needed to determine whether multiple \textit{wh-} fronting truly exists in the grammar as a productive multiple question formation strategy. The judgments reported in \REF{ex:matte:22} were observed stably from late 2022 to summer 2023, leading to the initial conclusion that multiple \textit{wh-} fronting exists in the language. However, in fall 2023, structures like those in \REF{ex:matte:22b} were consistently judged ungrammatical, leading to our present reluctance to fully commit to the existence of multiple \textit{wh-} fronting in the grammar. The third author of this paper (our language consultant) speculates that the availability of multiple \textit{wh-} fronting may have an inter-generational dimension, reporting that older Ekhwa Adara speakers like his mother or grandparents seemingly allow/produce multiple \textit{wh-} fronting structures. If all three multiple \textit{wh-} question formation strategies are indeed available in the grammar, strategies (i) and (ii) are clearly preferred over strategy (iii). Of relevance to the question of whether multiple \textit{wh-} fronting truly exists, we note that multiple focus movement does not seem to be available in declarative sentences, as shown in \REF{ex:matte:23}.

\ea \label{ex:matte:23}
\begin{xlist}

\ex \label{ex:matte:23a}
\gll * ɔnʃó\textsubscript{i} ku ijâ\textsubscript{j} ŋu omúsé sú mani {\longrule}\textsubscript{j} {\longrule}\textsubscript{i}\\
{} song	\Foc{} Ija \Foc{} Omuse \Fut{} teach {} {}\\
\glt Intended: `It is a \textsc{song} that Omuse will teach \textsc{ija}.' 

\ex \label{ex:matte:23b}
\gll * ijâ\textsubscript{i} ŋu/ku ɔnʃó\textsubscript{j} ku omúsé sú mani {\longrule}\textsubscript{i} {\longrule}\textsubscript{j}\\
{} Ija \Foc{} song \Foc{} Omuse \Fut{} teach {} {}\\
\glt Intended: `It is a \textsc{song} that Omuse will teach \textsc{ija}.'

\end{xlist}
\z

\noindent In order to emphasize more than one non-interrogative argument in the clause, the preferred strategy is to leave one item in-situ with a change in intonation (represented by all caps), as shown in \REF{ex:matte:24}.\footnote{A reviewer noted that the change in intonation in \REF{ex:matte:24} may serve a different focus function besides the emphasis of a non-interrogative argument as documented here (for example, marking an exhaustive argument). We acknowledge this point and leave the exploration of this possibility for future research.} 

\ea \label{ex:matte:24}
\gll ɔnʃó ku omúsé sú mani IJÂ\\
song \Foc{} Omuse \Fut{} teach Ija\\
\glt `It is a \textsc{song} that Omuse will teach \textsc{ija}.'
\z

\noindent The absence of multiple focus-fronting in this context potentially casts doubt on the existence of multiple \textit{wh-} fronting in the language (or at least in the grammar of the third author/others of his generation). If it does exist, something we leave for future research, then Ekhwa Adara is exceptional in its behavior regarding multiple \textit{wh-} questions, as it employs a multiple fronting strategy that is not found, as far as we know, in any other African language. 

\section{Conclusion} \label{sec:matte:4}

We summarize our key findings on question formation in Ekhwa Adara as follows. Polar questions in Ekhwa Adara are primarily formed using final-vowel lengthening and L\% boundary tone insertion. Furthermore, they are optionally marked by the sentence-initial Q particle \textit{kó}. The only exception to this rule is found in embedded polar questions, which require the presence of an embedded left peripheral Q particle and the absence of final-vowel lengthening + L\% boundary tones. Regarding movement, Ekhwa Adara is an optional \textit{wh-} movement language. When there is movement, the \textit{wh-} item obligatorily precedes a focus marker, which can also optionally be used with in-situ \textit{wh-} phrases. There is a selectional dependency between the \textit{wh-} item and the accompanying focus marker in all contexts. Ekhwa Adara allows for long-distance \textit{wh-} movement as well as simple (naked) partial \textit{wh-} movement. Indirect questions in the language are formed via relativization, except when embedded under the verb `ask', which involves \textit{wh-} movement to the left periphery of the embedded clause. Ekhwa Adara potentially has three distinct strategies for forming multiple \textit{wh-} questions: (i) all \textit{wh-} items remain in-situ; (ii) one \textit{wh-} item is fronted to clause-initial position, while the other(s) remain(s) in-situ; and (iii) multiple \textit{wh-} fronting, a phenomenon that to the best of our knowledge is (currently) unattested in any other African language. Multiple \textit{wh-} movement, to the extent that it truly exists in the grammar, is constrained – the initial fronted \textit{wh-} item must be hierarchically inferior to the other(s), an effect we have characterized as an anti-superiority effect. Despite the existence of the anti-superiority effect in multiple \textit{wh-} fronting constructions, \textit{wh-} movement in Ekhwa Adara is not constrained by the superiority condition that limits \textit{wh-} movement in languages like English.

Our investigation into Ekhwa Adara question syntax is ongoing. Future work will address the following issues (among others) related to the content of this paper: a fuller description of the distribution of optional Q particle \textit{kó}; a better understanding of the various focus particles and the selectional dependency between them and \textit{wh-} items; and an assessment of whether multiple \textit{wh-} fronting truly exists in the grammar and if so, a structural analysis of multiple \textit{wh-} fronting that derives the fact that the initial fronted \textit{wh-} item must be hierarchically inferior to the other(s), given the absence of superiority effects in the language. We hope our description of the Ekhwa Adara interrogative system inspires further research on the grammar of this considerably under-documented language.

\section*{Abbreviations}

\begin{tabularx}{.45\textwidth}{lQ}
\textsc{1} & first person\\
\textsc{3} & third person\\
\textsc{acc} & accusative\\
\textsc{comp} & complementizer\\
\textsc{def} & definite\\
\textsc{foc} & focus\\
\textsc{fut} & future\\
\end{tabularx}
\begin{tabularx}{.45\textwidth}{lQ}
\textsc{nom} & nominative\\
\textsc{pfv} & perfective\\
\textsc{prog} & progressive\\
\textsc{q} & question morpheme\\
\textsc{rel} & relative marker\\
\textsc{sg} & singular\\
\textsc{sp} & speaker
\end{tabularx}
\section*{Acknowledgements}

We thank our colleagues and the audience of ACAL 54 for helpful feedback, in particular, Daniel Aremu, Michael Erlewine, John Gluckman, Rebecca Jarvis, Travis Major, Malte Zimmerman, and two ACAL 54 proceedings reviewers. All authors contributed to the research reported in this article. Authors 1 and 2 wrote and presented the paper at ACAL 54 with assistance from author 10. Author 3 provided native speaker judgements of the data. All authors contributed to the collection and analysis of the data. Author 10 designed and managed the project.

\sloppy %this command eliminates a quirk that appears in the bibliography, you can just leave it here.

\printbibliography[heading=subbibliography,notkeyword=this]

\end{document}

